% Options for packages loaded elsewhere
\PassOptionsToPackage{unicode}{hyperref}
\PassOptionsToPackage{hyphens}{url}
\PassOptionsToPackage{dvipsnames,svgnames,x11names}{xcolor}
\PassOptionsToPackage{space}{xeCJK}
%
\documentclass[
  11pt,
]{article}
\usepackage{amsmath,amssymb}
\usepackage{iftex}
\ifPDFTeX
  \usepackage[T1]{fontenc}
  \usepackage[utf8]{inputenc}
  \usepackage{textcomp} % provide euro and other symbols
\else % if luatex or xetex
  \usepackage{unicode-math} % this also loads fontspec
  \defaultfontfeatures{Scale=MatchLowercase}
  \defaultfontfeatures[\rmfamily]{Ligatures=TeX,Scale=1}
\fi
\usepackage{lmodern}
\ifPDFTeX\else
  % xetex/luatex font selection
  \setmainfont[]{TeX Gyre Pagella}
  \setsansfont[]{TeX Gyre Heros}
  \setmonofont[]{Latin Modern Mono}
  \setmathfont[]{TeX Gyre Pagella Math}
  \ifXeTeX
    \usepackage{xeCJK}
    \setCJKmainfont[]{Noto Serif CJK SC}
          \setCJKsansfont[]{Noto Sans CJK SC}
              \setCJKmonofont[]{Noto Sans Mono CJK SC}
      \fi
  \ifLuaTeX
    \usepackage[]{luatexja-fontspec}
    \setmainjfont[]{Noto Serif CJK SC}
  \fi
\fi
% Use upquote if available, for straight quotes in verbatim environments
\IfFileExists{upquote.sty}{\usepackage{upquote}}{}
\IfFileExists{microtype.sty}{% use microtype if available
  \usepackage[]{microtype}
  \UseMicrotypeSet[protrusion]{basicmath} % disable protrusion for tt fonts
}{}
\makeatletter
\@ifundefined{KOMAClassName}{% if non-KOMA class
  \IfFileExists{parskip.sty}{%
    \usepackage{parskip}
  }{% else
    \setlength{\parindent}{0pt}
    \setlength{\parskip}{6pt plus 2pt minus 1pt}}
}{% if KOMA class
  \KOMAoptions{parskip=half}}
\makeatother
\usepackage{xcolor}
\usepackage[margin=1in]{geometry}
\usepackage{listings}
\newcommand{\passthrough}[1]{#1}
\lstset{defaultdialect=[5.3]Lua}
\lstset{defaultdialect=[x86masm]Assembler}
\usepackage{longtable,booktabs,array}
\usepackage{calc} % for calculating minipage widths
% Correct order of tables after \paragraph or \subparagraph
\usepackage{etoolbox}
\makeatletter
\patchcmd\longtable{\par}{\if@noskipsec\mbox{}\fi\par}{}{}
\makeatother
% Allow footnotes in longtable head/foot
\IfFileExists{footnotehyper.sty}{\usepackage{footnotehyper}}{\usepackage{footnote}}
\makesavenoteenv{longtable}
\setlength{\emergencystretch}{3em} % prevent overfull lines
\providecommand{\tightlist}{%
  \setlength{\itemsep}{0pt}\setlength{\parskip}{0pt}}
\setcounter{secnumdepth}{-\maxdimen} % remove section numbering
\ifLuaTeX
\usepackage[bidi=basic]{babel}
\else
\usepackage[bidi=default]{babel}
\fi
\babelprovide[main,import]{american}
\ifPDFTeX
\else
\babelfont{rm}[]{TeX Gyre Pagella}
\fi
% get rid of language-specific shorthands (see #6817):
\let\LanguageShortHands\languageshorthands
\def\languageshorthands#1{}
\usepackage{amsthm}
\usepackage{booktabs}
\usepackage{array}
\usepackage{tabularx}
\usepackage{longtable}
\usepackage{microtype}
\usepackage{fancyhdr}
\usepackage{lineno}
\usepackage{draftwatermark}
\usepackage[most]{tcolorbox}
\usepackage{listings}
\usepackage{xurl}
\usepackage{seqsplit}
\usepackage{enumitem}

\definecolor{OrbitBlue}{HTML}{163A5F}
\definecolor{OrbitLight}{HTML}{EDF4FA}
\definecolor{OrbitRule}{HTML}{8AA9C2}
\definecolor{OrbitCode}{HTML}{F6F7F9}
\definecolor{OrbitRed}{HTML}{9B2335}

\AtBeginDocument{\hypersetup{
  linkcolor=OrbitBlue,
  citecolor=OrbitBlue,
  urlcolor=OrbitBlue,
  pdfsubject={Finite algebra and circuit complexity review preprint},
  pdfkeywords={finite algebra, conservative operation, discriminator, term DAG, Shannon complexity, circuit depth}
}}

\SetWatermarkText{REVIEW DRAFT}
\SetWatermarkScale{0.55}
\SetWatermarkColor[gray]{0.95}
\SetWatermarkAngle{55}

\pagestyle{fancy}
\fancyhf{}
\fancyhead[L]{\small Fixed-$Q$ conservative term complexity}
\fancyhead[R]{\small Review draft}
\fancyfoot[C]{\thepage}
\setlength{\headheight}{14pt}

\newtcolorbox{reviewnotice}{
  colback=OrbitLight,
  colframe=OrbitRule,
  boxrule=0.6pt,
  arc=1.5pt,
  left=8pt,
  right=8pt,
  top=6pt,
  bottom=6pt
}

\newtheorem{theorem}{Theorem}[section]
\newtheorem{lemma}[theorem]{Lemma}
\newtheorem{proposition}[theorem]{Proposition}
\newtheorem{corollary}[theorem]{Corollary}
\theoremstyle{definition}
\newtheorem{definition}[theorem]{Definition}
\theoremstyle{remark}
\newtheorem{remark}[theorem]{Remark}

\lstset{
  basicstyle=\ttfamily\small,
  backgroundcolor=\color{OrbitCode},
  frame=single,
  rulecolor=\color{OrbitRule},
  breaklines=true,
  columns=fullflexible,
  keepspaces=true,
  showstringspaces=false,
  aboveskip=0.8em,
  belowskip=0.8em
}

\renewcommand{\arraystretch}{1.12}
\setlength{\emergencystretch}{1.5em}
\setlist{nosep, leftmargin=1.8em}
\modulolinenumbers[5]
\linenumbers
\ifLuaTeX
  \usepackage{selnolig}  % disable illegal ligatures
\fi
\IfFileExists{bookmark.sty}{\usepackage{bookmark}}{\usepackage{hyperref}}
\IfFileExists{xurl.sty}{\usepackage{xurl}}{} % add URL line breaks if available
\urlstyle{same}
\hypersetup{
  pdftitle={Exact Census and Simultaneous Shannon Size--Depth Bounds for Conservative Terms over a Fixed Three-Element Algebra},
  pdfauthor={AUTHOR METADATA PENDING HUMAN APPROVAL},
  pdflang={en-US},
  pdfkeywords={finite algebra, conservative
operation, discriminator, term DAG, Shannon complexity, local
coding, multivalued logic, circuit depth},
  colorlinks=true,
  linkcolor={Maroon},
  filecolor={Maroon},
  citecolor={Blue},
  urlcolor={Blue},
  pdfcreator={LaTeX via pandoc}}

\title{Exact Census and Simultaneous Shannon Size--Depth Bounds for
Conservative Terms over a Fixed Three-Element Algebra}
\author{AUTHOR METADATA PENDING HUMAN APPROVAL}
\date{Review draft --- 13 August 2026}

\begin{document}
\maketitle
\begin{abstract}
Fix the three-element algebra

\(Q=(\{0,1,2\};d,u)\),

where \(d(x,y,z)=z\) when \(x=y\) and \(d(x,y,z)=x\) otherwise, and
where \(u(0)=1\), \(u(1)=0\), and \(u(2)=1\). We study the
parameter-free \(r\)-ary term operations of \(Q\) that are conservative.
Their number is exactly

\(|\mathrm{CT}_r(Q)| = 2^{5\mathbin{\cdot}2^{r-1}-5} \mathbin{\cdot} 3^{3^r-3\mathbin{\cdot}2^r+3}\).

For an operation \(f\), let \(C_r(f)\) be the minimum number of
original-signature operation nodes in a rooted term DAG computing \(f\),
with free fan-out, and let \(D_r(f)\) be the corresponding minimum
operation depth. We give a parameter-free Lupanov-style compiler and
prove

\(\max_{f \in \mathrm{CT}_r(Q)} C_r(f) = \Theta(3^r/r)\).

The lower bound holds for almost every uniformly selected member of
\(\mathrm{CT}_r(Q)\). A sequence of depth refinements includes an exact
recursive signed-router family and an exact sibling-shared
program-vector construction. For \(q=3^w\), either fixed signed-router
vector has local size \(S(w)\) and depth \(D(w)\) satisfying

\(3\mathbin{\cdot}S(w) \le  4q+15\mathbin{\cdot}3^{\lceil w/2 \rceil}\)

and

\(D(w) \le  3+\lceil \log_2 w \rceil\).

Its leading size \((4/3+o(1))q\) matches a \(4q/3\) scalar-output lower
bound in the declared address-only model. This is an exact local
theorem, not a global compiler lower bound. An explicit two-plane
construction and a complement-relative binary construction culminate in
one and the same DAG having

\(\operatorname{size}=O(3^r/r)\)

and

\(\operatorname{depth} \le  r+O(\log r)\).

More explicitly, for \(r\ge 64\) the construction has size below
\(34\mathbin{\cdot}3^r/r\) and depth at most
\(r+4\mathbin{\cdot}\lceil \log_2 r \rceil+9\).

Writing \(\Delta_r(Q)=\max_{f \in \mathrm{CT}_r(Q)} D_r(f)\), syntax
counting also gives the almost-all lower bound

\(D_r(f) \ge  r-\log_3(\log r)-O(1)\),

and therefore the worst-case interval

\(r-\log_3(\log r)-O(1) \le  \Delta_r(Q) \le  r+O(\log r)\).

The displayed upper bound is a manuscript theorem with an explicit
generic proof and a separate no-author-import executable reconstruction.
It remains unformalized as an integrated Lean DAG theorem and does not
settle the optimal additive term. The independently audited
growing-router schedule supplied the earlier \(r+O(\sqrt{r})\) stage
after a required residual-first repair. Historical R6/R7 and R9 bounds
remain visible as prior compiler stages. The construction combines an
exact conservative-clone census, dynamic value names, classical local
coding, recursive subcube libraries, a balanced absorbing anchor, signed
discriminator routers, and parallel address programs. Deterministic
experiments check finite semantics, compiler schedules, mutations, and
structural arithmetic. They are validation evidence, not substitutes for
the generic proofs. No publication novelty, legal conclusion, practical
performance, or relationship to private Tau work is claimed.
\end{abstract}

{
\hypersetup{linkcolor=}
\setcounter{tocdepth}{2}
\tableofcontents
}
\begin{reviewnotice}
\textbf{Review and authorship status.} This PDF is a non-archival review copy.
Creator names, order, affiliations, corresponding-author details, and ORCIDs
remain subject to human approval. Do not deposit or represent this file as an
authored preprint until that gate is closed. The mathematical source is bound to
SHA-256 \texttt{01ddf4565212f798d0adf5c4c9b16d6de9074a317a34f2a39a80051de55a5410} and repository commit
\href{https://github.com/TheDarkLightX/OrbitSynthesis/commit/ab20b9df50f331e430c1a43a9d8c3af9c1e8556c}
{\texttt{ab20b9df50f3}}. The manuscript has not been externally peer
reviewed. Novelty, patent freedom to operate, license interpretation, and
practical performance remain unknown and are not claimed.
\end{reviewnotice}

\noindent\textbf{Keywords:} finite algebra; conservative operation;
discriminator; term DAG; Shannon complexity; local coding; multivalued
logic; circuit depth.

\noindent\textbf{MSC 2020:} 08A70 (applications of universal algebra in
computer science); 68Q06 (networks and circuits as models of
computation; circuit complexity).

\hypertarget{claim-status-at-a-glance}{%
\subsection*{Claim status at a glance}\label{claim-status-at-a-glance}}

\begin{longtable}[]{@{}
  >{\raggedright\arraybackslash}p{(\columnwidth - 4\tabcolsep) * \real{0.3333}}
  >{\raggedright\arraybackslash}p{(\columnwidth - 4\tabcolsep) * \real{0.3333}}
  >{\raggedright\arraybackslash}p{(\columnwidth - 4\tabcolsep) * \real{0.3333}}@{}}
\toprule\noalign{}
\begin{minipage}[b]{\linewidth}\raggedright
Object
\end{minipage} & \begin{minipage}[b]{\linewidth}\raggedright
Status in this draft
\end{minipage} & \begin{minipage}[b]{\linewidth}\raggedright
What is not claimed
\end{minipage} \\
\midrule\noalign{}
\endhead
\bottomrule\noalign{}
\endlastfoot
Exact census, given the stated Pixley specialization & unconditional
manuscript theorem & new discriminator interpolation theory \\
\(L_r(Q)=\Theta(3^r/r)\) for shared \(\{d,u\}\) DAGs & unconditional
manuscript theorem & an exact leading size constant or an ordinary-tree
theorem \\
\(\Delta_r(Q)=\Theta(r)\) and leading depth coefficient one &
unconditional manuscript theorem & an optimal additive-depth term \\
d-only signed family \(P_h,N_h\) & unconditional exact local theorem &
an optimal router in every grammar \\
sibling-shared program vector, Theorem 6.9 & unconditional exact local
theorem in its declared \(A=2\), scalar-output, free-fanout model &
simultaneous P/N cost, formula cost, or an integrated compiler lower
bound \\
\(O(3^r/r)\) size together with \(r+O(\log r)\) depth & unconditional
manuscript theorem; independently reconstructed, not end-to-end
Lean-formalized & an optimal additive-depth theorem, formula bound, or
bounded-fanout bound \\
novelty, FTO, patents, Tau/license relationship, practical speed &
\textbf{UNKNOWN / not claimed} & legal or commercial clearance \\
\end{longtable}

\hypertarget{introduction}{%
\subsection{1. Introduction}\label{introduction}}

Shannon complexity asks how large a circuit must be in the worst case,
and typically for almost every function, when the function class and
basis are fixed. Classical work treats Boolean circuits, \(k\)-valued
functional constructions, functionally complete bases, and local-coding
synthesis. The object here differs in two linked respects:

\begin{enumerate}
\def\labelenumi{\arabic{enumi}.}
\tightlist
\item
  the original signature \(\{d,u\}\) is parameter-free and preserves the
  proper subalgebra \(\{0,1\}\), so it is not functionally complete on
  \(Q\); and
\item
  the semantic class is the conservative part of the term clone of one
  fixed algebra.
\end{enumerate}

The first issue prevents us from simply naming constants. The second
makes the semantic entropy slightly smaller than that of all ternary
functions while retaining the leading term \((\log 3)3^r\). The
parameter-free obstacle is handled constructively: a nonbinary input
supplies dynamic names for \(0,1,2\), first locally and later through
one balanced global anchor.

The contribution stack established in this draft is:

\begin{enumerate}
\def\labelenumi{\arabic{enumi}.}
\tightlist
\item
  an exact closed-form census of the conservative \(r\)-ary term
  operations;
\item
  a complete coordinate-selector representation of that fragment;
\item
  worst-case and almost-all operation-node lower bounds of order
  \(3^r/r\);
\item
  a matching parameter-free shared-DAG compiler obtained by specializing
  classical Lupanov local coding;
\item
  worst-case and almost-all depth lower bounds of order \(r\);
\item
  an explicit compiler progression from exponential Gray dependency
  depth to linear depth, then to \(4r\), \(3r+O(\log r)\), the
  historical R9 bound \((3/2)r+O(r/\log r)\), the signed-router bound
  \(r+O(\sqrt{r})\), and finally the integrated parallel-program bound
  \(r+O(\log r)\), always preserving \(O(3^r/r)\) size on the same DAG;
\item
  exact six-way ternary, seven-way Boolean, and nine-way Boolean
  projection gadgets, retained as historical stages;
\item
  an exact d-only signed family \(P_h,N_h\) with \(3^{h-1}\)
  programmable branches at dependency depth h, independently audited and
  Lean-checked;
\item
  an exact sibling-shared absorbing-rail address summary that, for
  either fixed sign, materializes all controls for a capacity-q signed
  router in \((4/3+o(1))q\) shared-DAG nodes and \(O(\log \log q)\)
  preprocessing depth, with a matching leading scalar-output lower bound
  in the declared local model; and
\item
  source-auditable bounded checks with mutation controls, replay hashes,
  two independent program-vector specifications, and an independent
  reconstruction of the integrated compiler.
\end{enumerate}

The local-coding principle, recursive Shannon expansion, discriminator
interpolation, and general Shannon-counting methodology are not
contributions of this work. The paper-worthy candidate, subject to full
prior-art review, is their exact specialization and simultaneous
realization for this fixed parameter-free incomplete clone.

\hypertarget{algebra-semantic-class-and-cost-model}{%
\subsection{2. Algebra, semantic class, and cost
model}\label{algebra-semantic-class-and-cost-model}}

\hypertarget{the-fixed-algebra}{%
\subsubsection{2.1 The fixed algebra}\label{the-fixed-algebra}}

Let

\(Q=(\{0,1,2\};d,u)\)

with

\begin{lstlisting}
d(x,y,z) = z,  if x=y,
           x,  otherwise,

u(0)=1,  u(1)=0,  u(2)=1.
\end{lstlisting}

The binary set \(B=\{0,1\}\) is a proper subalgebra. Its nontrivial
automorphism is complement, written \(\overline{0}=1\) and
\(\overline{1}=0\). Notice that \(u\) is not conservative because
\(u(2)=1\).

An \(r\)-ary operation \(f:Q^r\to Q\) is \textbf{conservative} if

\(f(x_0,\ldots ,x_{r-1}) \in \{x_0,\ldots ,x_{r-1}\}\)

for every input tuple. Let \(\mathrm{CT}_r(Q)\) be the set of
conservative \(r\)-ary operations induced by parameter-free terms over
\(\{d,u\}\).

\hypertarget{shared-term-dags}{%
\subsubsection{2.2 Shared term DAGs}\label{shared-term-dags}}

A rooted original-signature term DAG has:

\begin{itemize}
\tightlist
\item
  \(r\) designated input nodes \(x_0,\ldots ,x_{r-1}\);
\item
  operation nodes labelled by unary \(u\) or ternary \(d\);
\item
  ordered incoming edges matching the operation arity; and
\item
  one designated output node.
\end{itemize}

Fan-out is free. Input nodes are not charged. For
\(f \in \mathrm{CT}_r(Q)\), define:

\begin{itemize}
\tightlist
\item
  \(C_r(f)\): the minimum number of operation nodes in such a DAG
  computing \(f\);
\item
  \(D_r(f)\): the minimum operation depth, with every input at depth
  zero and every operation one plus the maximum depth of its arguments;
\item
  \(L_r(Q)=\max_{f \in \mathrm{CT}_r(Q)} C_r(f)\); and
\item
  \(\Delta_r(Q)=\max_{f \in \mathrm{CT}_r(Q)} D_r(f)\).
\end{itemize}

The executable compiler reports all distinct reachable nodes, including
the \(r\) inputs, and puts inputs at depth one. For any fixed emitted
DAG, its total node count is its operation-node count plus at most
\(r\), and its reported depth is its analytic operation depth plus one.
The emitted values therefore upper bound, rather than equal, the minima
\(C_r(f)\) and \(D_r(f)\). All asymptotic theorems and displayed
analytic depth bounds use the operation-node convention.

Ordinary unshared term trees are not the main cost model. The baseline
compiler also has an expanded-tree analysis, but no order-optimal
ordinary-tree theorem is claimed here.

\hypertarget{semantic-characterization-and-exact-enumeration}{%
\subsection{3. Semantic characterization and exact
enumeration}\label{semantic-characterization-and-exact-enumeration}}

\hypertarget{pixley-specialization}{%
\subsubsection{3.1 Pixley specialization}\label{pixley-specialization}}

Pixley's internal-isomorphism characterization of quasi-primal term
operations {[}1,2{]} implies the following specialization for the fixed
algebra.

\textbf{Proposition 3.1 (conservative term criterion).} For \(r\ge 2\),
a conservative operation \(f:Q^r\to Q\) belongs to \(\mathrm{CT}_r(Q)\)
exactly when its restriction to the binary cube commutes with
simultaneous complement:

\(f(\overline{x_0},\ldots ,\overline{x_{r-1}}) = \overline{f(x_0,\ldots ,x_{r-1})}\)

for every \(x \in B^r\).

\textbf{Proof.} We spell out the internal-isomorphism catalog needed
from Pixley's criterion. The only nonempty proper subalgebra is
\(B=\{0,1\}\): it is closed under \(d\) and \(u\); a set containing
\(2\) and closed under \(u\) also contains \(1\) and then \(0\), so it
generates \(Q\); and neither singleton is closed under \(u\). Because
\(d\) is the discriminator, every permutation preserves \(d\). On \(B\),
the permutations preserving \(u\) are the identity and complement. On
\(Q\), the identity is the only permutation preserving \(u\): the unique
element \(2\) outside the two-cycle \(\{0,1\}\) must be fixed, after
which complement would send \(u(2)=1\) to \(0\) while \(u(2)\) remains
\(1\). There is no isomorphism between \(B\) and \(Q\) because their
cardinalities differ.

Thus the only nonidentity internal isomorphism that can constrain a term
operation is simultaneous complement on \(B^r\). Pixley's
characterization says that an operation preserving the relevant
subalgebras and internal isomorphisms is a term operation.
Conservativity supplies preservation of each tuple's value set and hence
of both \(B\) and \(Q\); on the binary cube the remaining condition is
exactly the displayed equivariance. Conversely every term operation
preserves subalgebras and commutes with their automorphisms, so the
condition is necessary. For \(r=1\), conservativity forces the identity
projection. This is a concrete application of classical theory, not a
new interpolation theorem. \(\square\)

\hypertarget{exact-census}{%
\subsubsection{3.2 Exact census}\label{exact-census}}

\textbf{Theorem 3.2 (exact conservative term count).} For every
\(r\ge 1\),

\(|\mathrm{CT}_r(Q)| = 2^{5\mathbin{\cdot}2^{r-1}-5} \mathbin{\cdot} 3^{3^r-3\mathbin{\cdot}2^r+3}\).

\textbf{Proof.} Partition \(Q^r\) by the set of values appearing in the
tuple.

On the binary cube, the \(2^r\) tuples form \(2^{r-1}\) complement
pairs. The constant pair \(\{0^r,1^r\}\) is forced by conservativity and
equivariance. Every other pair contains both values; choosing the value
on one representative determines the complementary value on its partner.
This contributes

\(2^{2^{r-1}-1}\).

A nonconstant tuple over \(\{0,2\}\) has two conservative output
choices, as does a nonconstant tuple over \(\{1,2\}\). Each alphabet
supplies \(2^r-2\) such tuples. These independent choices contribute

\(2^{2\mathbin{\cdot}(2^r-2)}\).

The all-\(2\) tuple is forced. Inclusion--exclusion gives

\(3^r-3\mathbin{\cdot}2^r+3\)

tuples using all three values, and each has three conservative output
choices. Their contribution is

\(3^{3^r-3\mathbin{\cdot}2^r+3}\).

Multiplication and collection of the powers of two give the formula. The
\(r=1\) instance equals one. \(\square\)

The first values are

\[
\begin{aligned}
|\mathrm{CT}_1(Q)| = 1, \\
|\mathrm{CT}_2(Q)| = 32, \\
|\mathrm{CT}_3(Q)| = 23,887,872.
\end{aligned}
\]

In particular,

\(\log |\mathrm{CT}_r(Q)| = (\log 3)3^r + O(2^r)\).

\hypertarget{coordinate-selector-completeness}{%
\subsubsection{3.3 Coordinate-selector
completeness}\label{coordinate-selector-completeness}}

\textbf{Proposition 3.3.} Every member of \(\mathrm{CT}_r(Q)\) has a
complement-invariant coordinate-selector table.

\textbf{Proof.} For every tuple \(x\), choose an index \(j(x)\) such
that \(f(x)=x_{j(x)}\). On each binary complement pair choose the index
on one representative and reuse it on its complement. Proposition 3.1
guarantees that the reused index produces the required complementary
value. Nonbinary internal-isomorphism orbits are singletons. Conversely,
the compilers below realize every such compatible selector table.
\(\square\)

Selector tables need not be unique: a tuple can contain the desired
output in several coordinates. The statement is semantic completeness,
not a canonical encoding claim.

\hypertarget{counting-lower-bounds}{%
\subsection{4. Counting lower bounds}\label{counting-lower-bounds}}

\hypertarget{shared-dag-size}{%
\subsubsection{4.1 Shared-DAG size}\label{shared-dag-size}}

\textbf{Theorem 4.1 (worst-case and almost-all size lower bound).} There
is a constant \(c>0\) such that

\(C_r(f) \ge  c\mathbin{\cdot}3^r/r\)

for a fraction tending to one of the operations
\(f \in \mathrm{CT}_r(Q)\). Consequently,

\(L_r(Q)=\Omega(3^r/r)\).

\textbf{Proof.} Topologically order a term DAG with \(s\) operation
nodes. At operation position \(i\), the number of choices for a unary
node and its predecessor is at most \(r+i\); the number for a ternary
node and its ordered predecessors is at most \((r+i)^3\). Choosing the
output and summing over all sizes at most \(s\) changes only constants,
so the number of descriptions is

\(\exp(O(s\mathbin{\cdot}\log(r+s)))\).

For \(s=c\mathbin{\cdot}3^r/r\), the exponent is
\(O(c\mathbin{\cdot}3^r)\). Theorem 3.2 gives
\(\exp((\log 3+o(1))\mathbin{\cdot}3^r)\) semantic operations. A
sufficiently small fixed \(c\) leaves an exponentially vanishing
fraction covered by the short descriptions. \(\square\)

This is a standard Shannon counting argument. It neither gives an
optimal leading constant nor an \(\Omega(3^r)\) lower bound.

\hypertarget{operation-depth}{%
\subsubsection{4.2 Operation depth}\label{operation-depth}}

Let \(A_h(r)\) overcount the number of original-signature term trees of
operation depth at most \(h\) on \(r\) variables. Unsharing a DAG does
not increase depth, and

\[
\begin{aligned}
A_0(r) = r, \\
A_{h+1}(r) \le  r + A_h(r) + A_h(r)^3 \le  3\mathbin{\cdot}A_h(r)^3
\end{aligned}
\]

for \(r\ge 2\). Therefore

\(\log A_h(r) \le  3^h\mathbin{\cdot}(\log r+(\log 3)/2)\).

\textbf{Theorem 4.2 (almost-all depth lower bound).} For a uniformly
selected \(f \in \mathrm{CT}_r(Q)\),

\(D_r(f) \ge  r-\log_3(\log r)-O(1)\)

with probability tending to one. In particular,

\(\Delta_r(Q)=\Omega(r)\).

\textbf{Proof.} If \(h\le r-\log_3(\log r)-K\) for a sufficiently large
fixed \(K\), the logarithm of the number of depth-\(h\) syntactic
candidates is a sufficiently small constant multiple of \(3^r\), whereas
Theorem 3.2 gives semantic log-cardinality \((\log 3+o(1))3^r\). The
fraction representable at that depth tends to zero exponentially in
\(3^r\). \(\square\)

The depth lower bound does not assume a bound on DAG size.

\hypertarget{parameter-free-local-coding}{%
\subsection{5. Parameter-free local
coding}\label{parameter-free-local-coding}}

\hypertarget{first-2-slices-and-dynamic-value-names}{%
\subsubsection{\texorpdfstring{5.1 First-\(2\) slices and dynamic value
names}{5.1 First-2 slices and dynamic value names}}\label{first-2-slices-and-dynamic-value-names}}

Partition the nonbinary tuples by the least coordinate \(a\) at which
\(x_a=2\). On this slice,

\begin{lstlisting}
x_0,...,x_(a-1) are binary,
x_a=2,
x_(a+1),...,x_(r-1) are unrestricted.
\end{lstlisting}

The slice has

\(N_a=2^a\mathbin{\cdot}3^{r-a-1}\)

rows. Within it, the terms

\begin{lstlisting}
u(u(x_a)),  u(x_a),  x_a
\end{lstlisting}

name \(0,1,2\), respectively. These are branch-local dynamic names, not
parameter-free constant terms on all of \(Q^r\).

The reachable first-\(2\) slices exactly partition the \(3^r-2^r\)
nonbinary tuples because

\(\sum_{a=0}^{r-1} 2^a\mathbin{\cdot}3^{r-a-1} = 3^r-2^r\).

This already yields a no-wasted-leaf classifier. In the executable's
total-node convention, writing \(T_r^{\mathrm{base}}\) for its emitted
node count, it gives

\[
\begin{aligned}
T_1^{\mathrm{base}}=6, \\
T_r^{\mathrm{base}}=3\mathbin{\cdot}3^r-3\mathbin{\cdot}2^{r-1}+6r-4 < 4\mathbin{\cdot}3^r  \text{for} r\ge 2.
\end{aligned}
\]

The corresponding fully expanded emitted term has recorded syntax size

\((35\mathbin{\cdot}7^r-9\mathbin{\cdot}3^r-50)/12\)

and depth \(4r+1\) in the executable convention. These are properties of
that compiler, not exact minima.

The constant-size equality selector used below is explicit:

\[
\begin{aligned}
\operatorname{EqSel}(x,y;p,q)=d(d(x,y,p),d(x,y,q),q).
\end{aligned}
\]

If \(x=y\), the outer discriminator returns \(p\); if \(x\ne y\), both
inner discriminators return \(x\) and the outer one returns \(q\). Thus
\(\operatorname{EqSel}\) has three operation nodes and two dependency
levels.

\hypertarget{mixed-domain-local-coding-lemma}{%
\subsubsection{5.2 Mixed-domain local-coding
lemma}\label{mixed-domain-local-coding-lemma}}

\textbf{Lemma 5.1 (mixed binary/ternary local coding).} Let

\(D=A_1 x \ldots  x A_n\),

where each \(|A_i|\) is two or three, and set \(N=|D|\). In a circuit
language with names for \(0,1,2\) and the displayed
\(\operatorname{EqSel}\), every table \(F:D\to Q\) has a shared circuit
of size \(O(N/n)\), uniformly over the mixture of coordinate alphabets.

\textbf{Construction.} For sufficiently large \(n\), put
\(H=\lfloor \log_3(N/n) \rfloor\); the finitely many smaller cases are
absorbed into the uniform constant. Choose a block of coordinates whose
number \(M\) of assignments is maximal subject to \(M\le H\). Maximality
gives \(H/3<M\le H\), while \(N\ge 2^n\) implies \(M=\Theta(n)\). Only
\(O(\log M)\) block coordinates are needed.

Precompute the point indicators for the \(M\) block assignments. There
are \(3^M\) functions from these assignments to \(Q\). A reflected
ternary Gray order changes one table entry at a time, so one
constant-size conditional update creates each successive library member
from the preceding one. Thus the shared library costs \(O(3^M)\), and

\(3^M \le  3^H \le  N/n\).

The remaining coordinates have \(N/M=O(N/n)\) assignments. A mixed-radix
decision tree classifies them and lets each leaf point into the shared
block-function library. Point indicators, the library, and the prefix
tree together cost \(O(N/n)\). \(\square\)

This is a concrete specialization of Lupanov's classical local-coding
principle {[}6{]}, not a new general synthesis method.

\hypertarget{size-optimal-parameter-free-compilation}{%
\subsubsection{5.3 Size-optimal parameter-free
compilation}\label{size-optimal-parameter-free-compilation}}

Apply Lemma 5.1 separately on every first-\(2\) slice, using the
slice-local names \(u(u(x_a)),u(x_a),x_a\) from Section 5.1. No global
anchor has been constructed or needed at this stage. On slice \(a\),
apply the lemma to its \(n=r-1\) varying coordinates; the \(r=1\)
identity is the separate finite base case. Uniformly absorbing the
finitely many small \(n\), the total nonbinary cost is

\(\sum_{a=0}^{r-1} O(N_a/r) = O(3^r/r)\).

The all-binary branch is classified modulo complement in \(O(2^r)\)
nodes, and the first-\(2\) dispatcher costs \(O(r)\). Both are lower
order. Together with Theorem 4.1 this proves:

\textbf{Theorem 5.2 (Shannon complexity).}

\(L_r(Q)=\Theta(3^r/r)\).

Moreover, for some constant \(c>0\), a uniformly selected member of
\(\mathrm{CT}_r(Q)\) has complexity at least \(c\mathbin{\cdot}3^r/r\)
with probability tending to one, while the compiler supplies the
matching order for every member.

The first executable realization uses a Gray-scheduled library. Its size
has the desired order, but the Hamiltonian dependency chain gives poor
depth. The next sections remove that scheduling artifact.

\hypertarget{simultaneous-size-and-depth}{%
\subsection{6. Simultaneous size and
depth}\label{simultaneous-size-and-depth}}

\hypertarget{layered-prefix-library}{%
\subsubsection{6.1 Layered prefix
library}\label{layered-prefix-library}}

Order the \(M\) block assignments as \(y_1,\ldots ,y_M\), with disjoint
point indicators \(chi_1,\ldots ,chi_M\). For every word \(w \in Q^j\),
build \(g_w\) so that it equals \(w_i\) at \(y_i\) for \(i\le j\) and
equals zero at all unprocessed points. The zero extension reuses its
parent; the one and two extensions each add one selector. Hence the
exact selector census is

\(2\mathbin{\cdot}\sum_{j=0}^{M-1} 3^j = 3^M-1\),

the same as for the Gray library, but the dependency depth is only \(M\)
selector levels. Combining this library with the prefix decision tree
proves that one DAG can simultaneously have

\(\operatorname{size}=O(3^r/r)\) and \(\operatorname{depth}=O(r)\).

Together with Theorem 4.2, this already establishes
\(\Delta_r(Q)=\Theta(r)\).

\hypertarget{recursive-subcube-library-and-the-4r-bound}{%
\subsubsection{\texorpdfstring{6.2 Recursive subcube library and the
\(4r\)
bound}{6.2 Recursive subcube library and the 4r bound}}\label{recursive-subcube-library-and-the-4r-bound}}

The layered construction is sharpened by recursively indexing the
library by its actual mixed-radix block coordinates. If their alphabet
sizes are \(q_1,\ldots ,q_b \in \{2,3\}\), let

\(M_j=\prod_{i=j}^b q_i\).

At a binary level, one \(\operatorname{EqSel}\) combines each pair of
suffix functions. At a ternary level, two nested
\(\operatorname{EqSel}\) terms combine each triple. The exact number of
selector calls is

\(S(q_1,\ldots ,q_b)=\sum_{j=1}^b (q_j-1)\mathbin{\cdot}3^{M_j} < 3\mathbin{\cdot}3^M\).

The top level dominates because \(M_{j+1}\le M_j/2\). The library
remains \(O(3^M)\), while a full mixed-radix prefix tree with \(P\)
leaves uses exactly \(P-1\) weighted selector calls.

The explicit depth accounting on first-\(2\) slice \(a\) is

\(2 + 2a + 4(r-a-1) + 2(a+1) = 4r\).

The four terms are, respectively, the deepest slice-local dynamic name,
binary and ternary coordinate routing inside the slice, and the
first-\(2\) dispatcher. The all-binary branch has depth at most
\(4r-2\).

\textbf{Theorem 6.1 (same-DAG \(4r\) construction).} Every
\(f \in \mathrm{CT}_r(Q)\) has one parameter-free original-signature DAG
with

\(\operatorname{size}=O(3^r/r)\) and \(\operatorname{depth}\le 4r\).

\hypertarget{one-global-nonbinary-anchor}{%
\subsubsection{6.3 One global nonbinary
anchor}\label{one-global-nonbinary-anchor}}

Define

\(h(x,y)=d(x,u(u(x)),d(y,u(x),x))\).

\textbf{Lemma 6.2.} For all \(x,y \in Q\), \(h(x,y)=2\) exactly when
\(x=2\) or \(y=2\); otherwise \(h(x,y)=x\).

Use an ordered balanced binary tree whose leaves, from left to right,
are \(x_0,\ldots ,x_{r-1}\), and label every internal node by \(h\);
call its root \(A\). The leaf order is part of the construction. On the
binary cube each \(h\) returns its left argument, so the tree returns
its distinguished leftmost leaf \(x_0\); an arbitrary reassociation with
\(x_0\) elsewhere is not asserted. Then

\[
\begin{aligned}
A=2 exactly on nonbinary tuples, \\
A=x_0 on the binary cube, \\
\operatorname{size}(A)=O(r), \\
\operatorname{depth}(A)\le 3\mathbin{\cdot}\lceil \log_2 r \rceil.
\end{aligned}
\]

On the nonbinary branch, \(u(u(A)),u(A),A\) are global dynamic names for
\(0,1,2\). On the binary branch they deliberately are not global
constants. This is what keeps the construction parameter-free and
compatible with the preserved subalgebra \(B\).

An exhaustive closure of binary-arity term functions through operation
depth three has cumulative sizes \(2,4,20,270\); the absorber's target
semantics first occurs at depth three. This finite fact concerns this
absorber, not the optimal depth of the final global compiler.

\hypertarget{a-branch-depth-three-ternary-router}{%
\subsubsection{6.4 A branch-depth-three ternary
router}\label{a-branch-depth-three-ternary-router}}

On a branch where terms named \(0,1,2\) have the indicated values,
define

\begin{lstlisting}
mux(x,b0,b1,b2)
 = d(
     d(0,d(x,2,b0),d(0,x,b0)),
     0,
     d(d(1,x,b1),1,d(x,2,b2))
   ).
\end{lstlisting}

\textbf{Lemma 6.3.} For every \(x,b0,b1,b2 \in Q\),

\(mux(x,b0,b1,b2)=b_x\).

Every branch argument lies exactly three discriminator nodes below the
root. Direct substitution for \(x=0,1,2\) proves the identity. With
\(0=u(u(A))\), \(1=u(A)\), and \(2=A\), the displayed shared
materialization has seven \(d\) nodes and the two shared \(u\) name
nodes. It therefore has nine operation nodes, or fourteen total nodes if
its five distinct raw terminals are also counted. Its operation depth is
five relative to an input anchor. An older executable reported fifteen
because it constructed \(u(A)\) twice; that is an emitter artifact, not
the shared-DAG census used here.

A finite Z3 \passthrough{\lstinline!QF\_BV!} search reports
\passthrough{\lstinline!UNSAT!} for branch depth at most two in a
deliberately generous selector-only grammar: all 27 unary functions of
the selector are admitted as free terminals; \(b0,b1,b2\) are branch
terminals; internal nodes are \(u\) or \(d\); and all 81 valuations are
imposed on the complete depth-two syntax tree. Together with the
displayed term, this makes depth three optimal only for that finite
branch-scoped grammar. It is not a global depth lower bound and
currently has no separately checked proof certificate.

\hypertarget{global-library-and-final-depth-theorem}{%
\subsubsection{6.5 Global library and final depth
theorem}\label{global-library-and-final-depth-theorem}}

Put

\(H=\lfloor \log_3(3^r/r) \rfloor\).

Choose \(b\) ternary coordinates so that their number of assignments
\(M=3^b\) is maximal subject to \(M\le H\). Apart from finitely many
small arities, \(H/3<M\le H\), and therefore

\[
\begin{aligned}
M=\Theta(r), \\
3^M\le 3^r/r, \\
3^r/M=O(3^r/r).
\end{aligned}
\]

Build every \(Q\)-valued function on the local block recursively, using
the router of Lemma 6.3 at each coordinate. The library costs
\(O(3^M)\). Route the remaining prefix coordinates with the same router;
the prefix tree costs \(O(3^r/M)\). Add the \(O(r)\) global anchor and
the lower-order \(O(2^r)\) complement-orbit classifier for the binary
branch. A final \(\operatorname{EqSel}\) chooses that binary classifier
exactly when \(u(u(A))=A\) and otherwise chooses the global nonbinary
library.

The block and prefix partition the \(r\) routed coordinates. Each routed
coordinate adds three branch-dependency levels. Put
\(\alpha=3\mathbin{\cdot}\lceil \log_2 r \rceil\). The deepest dynamic
name has depth \(\alpha+2\), and the nonbinary routed output before the
final choice has depth at most \(3r+\alpha+2\). The binary
complement-orbit classifier uses at most \(r-1\) sequential equality
decisions, each one \(\operatorname{EqSel}\) of branch-dependency depth
two, so its depth is at most \(2(r-1)\). The final
\(\operatorname{EqSel}\) adds two levels. Therefore the second compiler
has depth

\begin{lstlisting}
2+max(alpha+2, 2(r-1), 3r+alpha+2)
  =3r+3*ceil(log_2 r)+4.
\end{lstlisting}

Taking the better of this construction and Theorem 6.1 gives:

\textbf{Theorem 6.4 (simultaneous Shannon size and explicit depth).}
Every \(f \in \mathrm{CT}_r(Q)\) has one parameter-free
original-signature term DAG with

\(\operatorname{size}=O(3^r/r)\)

and

\(\operatorname{depth}\le min(4r,3r+3\mathbin{\cdot}\lceil \log_2 r \rceil+4)\).

Combining Theorems 4.2 and 6.4 gives

\(r-\log_3(\log r)-O(1) \le  \Delta_r(Q) \le  min(4r,3r+3\mathbin{\cdot}\lceil \log_2 r \rceil+4)\),

and in particular

\(r-\log_3(\log r)-O(1) \le  \Delta_r(Q) \le  3r+O(\log r)\).

Thus this historical compiler already proves \(\Delta_r(Q)=\Theta(r)\)
with leading-coefficient interval \([1,3]\). The historical fixed-router
stages below strictly narrowed that upper endpoint before the
signed-family construction. Neither endpoint is claimed optimal.

\hypertarget{historical-fixed-router-stages}{%
\subsubsection{6.6--6.7 Historical fixed-router
stages}\label{historical-fixed-router-stages}}

Two fixed depth-three routers preceded the recursive signed family. They
are retained for provenance and comparison, not because the optimized
theorem depends on them.

\begin{itemize}
\tightlist
\item
  \(R6\) has 13 discriminator nodes and six ternary branches. Its six
  fixed 18-slot programs realize all branch projections.
\item
  \(R7\) has 12 discriminator nodes and seven Boolean branches. Its
  programs use the complement-relative binary names \(x_0,u(x_0)\), not
  absolute constants.
\item
  \(R9\) is a full 13-node depth-three discriminator tree with nine
  Boolean branches and 18 program slots. Nine programs realize all nine
  projections.
\end{itemize}

Balanced six-way routing gave the historical nonbinary ceiling

\((3\mathbin{\cdot}\log_6(3))r+O(r/\log r)=1.839441578296\ldots r+O(r/\log r)\).

For R9, encode the dynamic values by \(0\to 00,1\to 01,2\to 10\), run
two copies in parallel, and decode once with

\(\operatorname{Dec}(high,low)=d(d(high,one,two),zero,low)\).

This gave the later historical ceiling

\((3/2)r+O(r/\log r)\).

At every route level one shared address-control family serves all active
routers (and both R9 planes), so the fixed-router schedules preserve
\(O(3^r/r)\) size on the same DAG. The exact telescoping router census,
balanced-library recurrences, residual schedules, and all fixed programs
are in \path{notes/QUASIPRIMAL_CONSERVATIVE_TERM_PROGRAMMABLE_DEPTH.md}
and \path{notes/QUASIPRIMAL_CONSERVATIVE_TERM_BOOLEAN_R9_DEPTH.md}.
Deterministic replay covers 4,374 R6 cases, 896 R7 cases, 4,608 R9
projection cases, 9,216 complement-relative cases, and 177,147 two-plane
cases. These bounded checks validate the gadgets; the all-arity schedule
statements remain manuscript proofs and historical calibration.

The old bounds were, respectively,

\[
\begin{aligned}
\operatorname{size}=O(3^r/r),  \operatorname{depth}\le (3\mathbin{\cdot}\log_6(3))r+O(r/\log r), \\
\operatorname{size}=O(3^r/r),  \operatorname{depth}\le (3/2)r+O(r/\log r).
\end{aligned}
\]

Neither fixed router is claimed capacity-optimal. A radius-one search
found no R10 insertion into the frozen R9, which is only a bounded
\passthrough{\lstinline!NO\_HIT!}.

\hypertarget{recursive-signed-discriminator-routers}{%
\subsubsection{6.8 Recursive signed discriminator
routers}\label{recursive-signed-discriminator-routers}}

The fixed R9 is not the final routing object. On the Boolean subalgebra
define positive and negative depth-one cells

\[
\begin{aligned}
P_1(x;c0,c1)=d(x,c0,c1), \\
N_1(x;c0,c1)=d(c0,x,c1).
\end{aligned}
\]

The exact programs are

\begin{longtable}[]{@{}llll@{}}
\toprule\noalign{}
cell & projection mode & constant 0 & constant 1 \\
\midrule\noalign{}
\endhead
\bottomrule\noalign{}
\endlastfoot
\(P_1\) & \((0,0)\) returns x & \((1,0)\) & \((0,1)\) \\
\(N_1\) & \((0,1)\) returns \(1-x\) & \((0,0)\) & \((1,1)\) \\
\end{longtable}

Recursively set

\[
\begin{aligned}
P_h=d(P_{h-1},N_{h-1},P_{h-1}), \\
N_h=d(N_{h-1},P_{h-1},N_{h-1}).
\end{aligned}
\]

For a target in the left, middle, or right branch group, program the
three children respectively as

\begin{lstlisting}
(projection,constant 0,constant 0),
(constant 0,projection,constant 1),
(constant 0,constant 0,projection).
\end{lstlisting}

At a positive parent the resulting expressions are

\[
\begin{aligned}
d(x,0,0)=x, \\
d(0,1-x,1)=x, \\
d(0,0,x)=x.
\end{aligned}
\]

The same transitions at a negative parent return \(1-x\). Constants lift
by \(d(c,c,c)=c\).

\textbf{Theorem 6.7 (exact strong signed-router family).} For every
\(h\ge 1\), there are d-only full trees \(P_h,N_h\) with

\begin{lstlisting}
capacity                 q_h=3^(h-1),
branch leaves             q_h,
address-control leaves    2q_h,
discriminator nodes       (3^h-1)/2,
dependency depth          h.
\end{lstlisting}

For every target, \(P_h\) returns the target branch and \(N_h\) its
Boolean complement; both signs also have constant-zero and constant-one
modes. Every branch path in \(P_h\) has even middle-edge parity, while
every branch path in \(N_h\) has odd parity.

\textbf{Proof.} The semantic modes follow from the preceding induction.
The recurrences

\[
\begin{aligned}
q_h=3q_{h-1}, \\
C_h=3C_{h-1}, \\
S_h=1+3S_{h-1}, \\
D_h=1+D_{h-1}
\end{aligned}
\]

start at \(q_1=1,C_1=2,S_1=D_1=1\) and give the displayed counts. The
middle child swaps P and N, while the outer children preserve the sign,
proving the path-parity statement. No NOT or signed-leaf constructor
occurs in the tree. \(\square\)

Thus depth four already gives an exact R27 with 40 d nodes, 27 branches,
and 54 controls; depth six gives R243. At fixed h the routing
coefficient is \(h/(h-1)\), tending to one. The family theorem has been
independently derived and replayed and has also been formalized in Lean
in the exact d-only grammar. The grammar's lack of a constant or NOT
constructor is enforced by its inductive constructors, not by the
historical \passthrough{\lstinline!OriginalSignature!} predicate, which
is true of every inhabitant. A later additive Lean layer re-proves the
address/control cardinalities and \(q_h=3^{h-1}\) with kernel-reduced
\passthrough{\lstinline!decide!} and induction, avoiding the predecessor
file's native-evaluator checkpoints.

The representation bridges must still be charged. On a nonbinary slice,
\(A=2\) supplies legal names \(u(A)\) and \(u(u(A))\); two sibling
copies of the router carry the Boolean planes and are decoded once. On
the binary cube a logical relative program bit p is address-only, but
its physical representative is \(x_0\) when \(p=0\) and \(u(x_0)\) when
\(p=1\). Those physical wires are payload-relative. Their legality comes
from simultaneous-complement equivariance, not from treating them as
absolute constants or literally payload-independent controls.

Choosing \(k=\lfloor \sqrt{r} \rfloor\), \(h=k+1\), and \(q=3^k\), and
routing every short residual chunk first, gives the following audited
compiler implication.

\textbf{Conditional Corollary 6.8 (growing-router compiler).} Assume the
coordinate-selector, global-anchor, two-plane representation,
binary-orbit, decoder, and outer-glue lemmas used by the preceding
compiler. Then every \(f \in \mathrm{CT}_r(Q)\) has one
original-signature shared DAG with

\[
\begin{aligned}
\operatorname{size}  = O(3^r/r), \\
\operatorname{depth} = r+O(\sqrt{r}).
\end{aligned}
\]

The residual-first qualification is necessary for this ledger. A
residual-last schedule has an explicit coefficient-breaking
counterexample at \(r=512\). An independent accounting audit checks the
repaired integer schedule for every \(64\le r\le 10000\), including all
controls, both planes, padding, the binary branch, anchor, decoder, and
glue. This is an audited conditional composition, not a Lean proof of
the end-to-end compiler.

\hypertarget{sibling-shared-program-vectors-and-logarithmic-overhead}{%
\subsubsection{6.9 Sibling-shared program vectors and logarithmic
overhead}\label{sibling-shared-program-vectors-and-logarithmic-overhead}}

The signed family exposes a second bottleneck: compiling \(2q\) program
slots one at a time costs \(\Theta(q \log q)\) for a capacity-q router.
The slots instead form a classical first-nonpropagating prefix state.
The balanced-prefix idea is standard {[}29,30{]}; the local content here
is its charged constant-free \(\{d,u\}\) realization and exact sharing
census.

A \textbf{multi-output shared DAG} is an acyclic original-signature
operation graph with a finite ordered list of distinguished roots. Its
size is the number of operation nodes in the union of their ancestor
subgraphs, counted once; inputs are free and fan-out is unrestricted.
Substitution adjoins that union once and permits arbitrary references to
its roots. The final compiler still has one output root and uses the
single-output convention of Section 2.

Write a physical branch and target as \(p,t \in \{0,1,2\}^w\), and put
\(q=3^w\). Let \(B_p(t)\) be one when the first mismatch exists and is
not \((t_j,p_j)=(1,2)\), and let \(G_p(t)\) be one when the first
mismatch is exactly that exceptional pair. Equality is the residual
state \((B,G)=(0,0)\). The rails are disjoint. For consecutive word
blocks \(U,V\), first-mismatch ownership gives the associative laws

\[
\begin{aligned}
B_{\mathrm{UV}}=d(B_U,G_U,B_V), \\
G_{\mathrm{UV}}=d(G_U,B_U,G_V).
\end{aligned}
\]

Writing \(N=\neg B\), the complement rail co-produces in the same layer:

\[
\begin{aligned}
N_{\mathrm{UV}}=d(G_U,B_U,N_V).
\end{aligned}
\]

On the nonbinary branch put \(one=u(A)\) and \(zero=u(one)\). The
complete one-digit base uses four non-name operation nodes:

\[
\begin{aligned}
\text{physical} 0: B=d(x,A,one),       G=zero, \\
\text{physical} 1: B=d(one,x,zero),    G=zero, \\
\text{physical} 2: B=u(d(x,A,one)),    G=d(x,A,zero).
\end{aligned}
\]

The final control pair is \((B,G)\) in a positive cell and \((G,N)\) in
a negative cell. Cell sign is the root sign xor the parity of the middle
digits of \(p\). Constant mode \(c\) uses \((1-c,c)\) in a positive cell
and \((c,c)\) in a negative cell. No payload wire enters this address
program.

Two identities create the decisive shared-DAG reduction:

\[
\begin{aligned}
G_{a s}=G_a       \text{when} the \text{physical} suffix s contains no digit 2, \\
N_{a1}=G_{a2}.
\end{aligned}
\]

The first reuses the left gain root whenever the right gain rail is
zero. The second points every eligible not-B rail to an already
materialized sibling gain root. Both are ordinary root reuse under free
fan-out, not free gates.

Put \(a=\lceil w/2 \rceil\), \(b=\lfloor w/2 \rfloor\), and \(A_0=3^a\).
Let \(R(w)\) count the full \((B,G)\) library and \(E(w)\) the full
\((B,G,N)\) suffix library, excluding the two names. Exact structural
counting gives

\[
\begin{aligned}
R(1)=4, \\
R(w)=R(a)+R(b)+3^w+A_0\mathbin{\cdot}(3^b-1)/2, \\
E(1)=4, \\
E(w)=R(a)+E(b)+2\mathbin{\cdot}3^w-3^{w-1}+A_0\mathbin{\cdot}(3^b-1)/2.
\end{aligned}
\]

At the top merge, the mixed/gain overlap has size

\[
\begin{aligned}
O_P(w)=(3^{w-1}+1)/2, \\
O_N(w)=(3^{w-1}-1)/2.
\end{aligned}
\]

Thus, for \(w\ge 2\), the exact sizes of the displayed P and N
constructions are

\[
\begin{aligned}
S_P(w)=2+R(a)+E(b)+A_0\mathbin{\cdot}(3^b-1)/2+3^w-O_P(w), \\
S_N(w)=2+R(a)+E(b)+A_0\mathbin{\cdot}(3^b-1)/2+3^w-O_N(w).
\end{aligned}
\]

The direct bases are \(S_P(1)=5\) and \(S_N(1)=6\).

\textbf{Theorem 6.9 (sibling-shared program vector).} Let \(A=2\). For
either fixed P or N root sign and every \(w\ge 1\), all \(2q\)
projection controls of the capacity-\(q=3^w\) signed router have one
legal address-only scalar \(\{d,u\}\) multi-output shared DAG satisfying

\[
\begin{aligned}
3\mathbin{\cdot}S(w) \le  4q+15\mathbin{\cdot}3^{\lceil w/2 \rceil}, \\
S(w)   \le  7q/3, \\
D(w)   \le  3+\lceil \log_2 w \rceil,
\end{aligned}
\]

including the two shared name nodes and measured above raw
anchor/address inputs. Consequently \(S(w)=(4/3+o(1))q\). Width zero
uses only the two names.

\textbf{Proof.} Induction on the balanced split gives
\(R(w)\le 20\mathbin{\cdot}3^w/9\) and
\(E(w)\le 26\mathbin{\cdot}3^w/9\); after substituting these
inequalities in the exact top count, the residual half-width term is at
most \(5\mathbin{\cdot}3^{\lceil w/2 \rceil}\). This proves the first
display. It implies \(S\le 7q/3\) from width four onward, and widths one
through three are direct bases. The one-digit rails have depth at most
three, and each balanced merge adds one level, proving the depth bound.
\(\square\)

The bound has a matching leading lower term in this declared local
model. The gain family has \((q+1)/2\) distinct scalar functions, the
sign-selected B/N family has \(q\), and the only cross-family identities
are \(N_{a1}=G_{a2}\). Their counts are \((q/3+1)/2\) for P and
\((q/3-1)/2\) for N. Hence the requested controls comprise exactly
\(4q/3\) distinct functions for P and \(4q/3+1\) for N. Every one is
Boolean-valued, whereas the free terminals are the constant anchor
\(A=2\) and Q-surjective address coordinates. No requested control is a
free terminal, and one scalar operation node can root only one new
scalar function. Therefore the leading local size constant is exactly
\(4/3\).

This lower bound does not apply to an arbitrary integrated router that
avoids exposing these exact controls, to simultaneous P and N
materialization, to formulas, or to bounded fan-out. It does not
determine the finite-width exact minimum.

Two independent no-author-import reconstructions check the semantics,
recurrences, sharing identities, and lower census. Lean proves the rail
algebra, one-digit bases, zero-gain reuse, sibling identity, final
control decoding, exact arithmetic recurrences, the all-width envelope,
the uniform upper bound, and the matched sharp depth
\(3+\lceil \log_2 w \rceil\). The formal cost object is still a declared
recurrence ledger rather than a serialized hash-consed DAG; its
refinement link is the executable construction plus independent
structural audit. The distinct-output lower census is only partially
formalized: the bad/not-bad families are injective, but the gain
quotient and final signed inclusion-exclusion remain open in Lean. A
direct one-wire three-state encoding with merge \(d(S_U,A,S_V)\) is also
valid, but its per-cell decoding is strictly more expensive than the
sibling-shared two-rail construction.

To apply this vector to the compiler, fix \(r\ge 64\), put

\[
\begin{aligned}
L=4+\lceil \log_3(r^2) \rceil, \\
H=r-L, \\
M=3^{\lfloor \log_3 H \rfloor}, \\
b=\log_3 M, \\
s=r-b, \\
P=3^s,
\end{aligned}
\]

and let \(C=\lceil \log_2 r \rceil\). Proposition 3.3 supplies a total
coordinate selector \(\sigma:Q^r\to \{0,\ldots ,r-1\}\) that is
invariant under simultaneous complement on the Boolean cube. Lemma 6.2
supplies the ordered balanced anchor \(A\), with \(A=2\) off that cube
and \(A=x_0\) on it.

\textbf{Lemma 6.10a (nonbinary two-plane compiler).} On \(A=2\), there
is one legal original-signature shared DAG whose output is
\(x_{\sigma(x)}\). Inclusive of the anchor, two generated names, one
decoder, and the three nodes reserved for the final glue, its
operation-node count is at most

\[
\begin{aligned}
N_{\mathrm{NB}} = (3M-1)3^M+(3P-1)+S_P(b)+S_P(s)-2+4(r-1)+2+3,
\end{aligned}
\]

and hence is below \((32+1/27+1/2)3^r/r\). Its output before the final
glue has depth at most \(r+4C+7\).

\textbf{Proof.} Write each input as \(x=(p,l)\) with \(p \in Q^s\) and
\(l \in Q^b\). For a fixed prefix \(p\), define the local table

\[
\begin{aligned}
g_p(l)=(p,l)_{\sigma(p,l)}.
\end{aligned}
\]

Materialize the universal library of all \(3^M\) Q-valued tables on
\(Q^b\). Encode table values by the two legal Boolean planes

\[
\begin{aligned}
0 \to  (0,0),   1 \to  (0,1),   2 \to  (1,0).
\end{aligned}
\]

Each table uses two capacity-M P routers, controlled by \(l\); every
table and both planes share the same width-b vector from Theorem 6.9.
For each prefix \(p\), take the two roots belonging to \(g_p\). Two
capacity-P P routers, controlled by \(p\) and sharing one width-s
vector, select those roots. Decode once with

\[
\begin{aligned}
\operatorname{Dec}(high,low)=d(d(high,one,A),zero,low).
\end{aligned}
\]

Its values on \(00,01,10\) are respectively \(0,1,2\); the unused code
\(11\) is irrelevant. Thus the output is exactly
\(g_p(l)=x_{\sigma(x)}\). Every control is a term over the raw address
and \(A\); the names \(one=u(A)\) and \(zero=u(one)\) are charged nodes,
not nullaries. Each vector is adjoined once and reused by every table,
plane, and router that names it.

Two capacity-M routers have \(3M-1\) discriminator nodes, and two
capacity-P routers have \(3P-1\). This gives the displayed exact upper
ledger. Since \(\lceil \log_3(r^2) \rceil\le 2\log_3(r)+1\) and
\(5+2\log_3(r)<r/2\) from \(r=64\) onward, \(H>r/2\); together with
\(H/3<M\le H\), this gives \(M>r/6\). The two-logarithm reserve also
gives

\[
\begin{aligned}
3^M \le  3^r/(81r^2).
\end{aligned}
\]

Consequently the local routers cost less than \((1/27)3^r/r\), while the
prefix routers and width-s vector cost less than

\[
\begin{aligned}
(3+7/3)P < 32\mathbin{\cdot}3^r/r.
\end{aligned}
\]

The width-b vector and fixed overhead are below \(7r\) nodes. The
elementary inequality \(14r^2<3^r\), valid from \(r=64\), puts them
below \((1/2)3^r/r\).

The anchor has depth \(3C\). With
\(D_P(w)\le 3+\lceil \log_2 w \rceil\), the two plane depths before
decoding satisfy

\[
\begin{aligned}
D_{\mathrm{local}}  \le  3C+D_P(b)+b+1, \\
D_{\mathrm{prefix}} \le  \max(D_{\mathrm{local}},3C+D_P(s))+s+1.
\end{aligned}
\]

Adding the two decoder levels and using \(b+s=r\) gives
\(D_{\mathrm{NB}}\le r+4C+7\). \(\square\)

The binary cube cannot use the absolute names just used. Its compiler is
instead complement-relative.

\textbf{Lemma 6.10b (binary complement-relative compiler).} On the
Boolean cube there is a legal original-signature shared DAG whose output
is \(x_{\sigma(x)}\). For \(r\ge 64\) its increment over the already
present anchor is below \((1/2)3^r/r\), and its intrinsic output depth
is at most \(r+4C+7\).

\textbf{Proof.} Put \(z_i=0\) when \(x_i=x_0\) and \(z_i=1\) otherwise.
The word \(z\) and the selector index are invariant under simultaneous
complement. Choose

\[
\begin{aligned}
k=\lfloor sqrt r \rfloor,  q=3^k,  w=\lfloor \log_2 q \rfloor.
\end{aligned}
\]

Partition the \(r-1\) relative coordinates into chunks of width at most
\(w\), placing the short residual chunk first. At a chunk of width
\(c\), fix an injection \(\iota_c:\{0,1\}^c\to \{0,1,2\}^k\) into the
first \(2^c\) branches of a capacity-q P router, and point every unused
branch to the last live child. Each logical control table
\(f:\{0,1\}^c\to \{0,1\}\) is the corresponding Theorem 6.9 P-control
coordinate for \(\iota_c(z)\). It is installed as the physical,
complement-relative wire \(x_0 \mathbin{\mathrm{xor}} f(z)\) and
realized from leaves \(x_0,u(x_0)\) by the decision node

\[
\begin{aligned}
\operatorname{Sel}(x_i,x_0,E,D)=d(d(x_i,x_0,E),d(x_i,x_0,D),D),
\end{aligned}
\]

which returns \(E\) when \(x_i=x_0\) and \(D\) otherwise. One physical
control therefore costs at most \(3(2^c-1)\) nodes and depth \(2c+1\).
All router instances at a level share the same \(2q\) controls. On the
representative orientation \(x_0=0\), these are the ordinary program
bits. At the terminal prefix \(z\), use leaf \(x_{\sigma(0,z)}\); the
routed leaf is therefore the requested coordinate. Complementing the
input complements every payload and physical control; self-duality of
\(d\) on \(\{0,1\}\) and complement invariance of \(\sigma\) then prove
the other orientation. No absolute Boolean constant has been used.

If \(I\) is the total number of router instances and the chunk widths
are \(c_j\), residual-first ordering gives \(I<2^{r-1}\). Hence

\[
\begin{aligned}
B_{\mathrm{router}}  = ((3q-1)/2)I                 < (1/4)3^r/r, \\
B_{\mathrm{control}} = \sum_j 6q(2^{c_j}-1)         < (1/4)3^r/r.
\end{aligned}
\]

For completeness, the first inequality follows from \(3rq2^r<3^r\); the
second follows from \(B_{\mathrm{control}}<6rq^2\) and \(24r^2q^2<3^r\).
Since \(q\le 3^{\sqrt{r}}\), these reduce to positivity of

\[
\begin{aligned}
F_1(r)=(1-\log_3 2)r-\sqrt{r}-\log_3(3r), \\
F_2(r)=r-2\sqrt{r}-\log_3(24r^2).
\end{aligned}
\]

Both are positive at 64. For \(r\ge 64\),

\[
\begin{aligned}
F_1'(r)=1-\log_3 2-1/(2\sqrt{r})-1/(r ln 3)>0, \\
F_2'(r)=1-1/\sqrt{r}-2/(r ln 3)>0,
\end{aligned}
\]

so the bounds hold for every real \(r\ge 64\), hence every integer arity
in scope. The sole \(u(x_0)\) leaf is an existing anchor subnode and is
not charged twice.

If \(m\) is the number of chunks, the intrinsic depth is at most

\[
\begin{aligned}
D_B \le  (k+1)m+2w+1.
\end{aligned}
\]

For \(k\ge 8\), \(3k/2\le w<8k/5\). Indeed, \(2^19<3^12\) and
\(3^5<2^8\) give \(19/12<\log_2(3)<8/5\), and the lower fractional
margin exceeds one half already at \(k=8\). Since \(m\le 2(r-1)/(3k)+1\)
and \(r<(k+1)^2\), this yields

\[
\begin{aligned}
D_B < 2r/3+(73/15)k+8/3 \\
\le  2r/3+(73/15)\sqrt{r}+8/3 \\
\le  r+4C+7.
\end{aligned}
\]

For the last inequality, use \(C\ge \log_2 r\). The difference is at
least

\begin{lstlisting}
r/3+4log_2(r)+13/3-(73/15)sqrt(r).
\end{lstlisting}

It is positive at 64, and its derivative is at least
\(1/3-73/(30\sqrt{r})>0\) there and thereafter. The displayed upper
ledger is \(\log_3(2)r+O(sqrt r)\), so this branch is asymptotically
below \(r\). \(\square\)

Put \(zero=u(u(A))\), let \(B\) and \(N\) denote the binary and
nonbinary outputs, and define

\[
\begin{aligned}
\operatorname{Glue}(A,B,N)=d(d(zero,A,B),d(zero,A,N),N).
\end{aligned}
\]

The glue adds exactly three discriminator nodes and two levels. On the
Boolean cube \(zero=A\), so it returns \(B\); off the cube
\(A=2,zero=0\), so it returns \(N\).

\textbf{Theorem 6.10 (integrated parallel-program compiler).} Every
\(f \in \mathrm{CT}_r(Q)\) has one parameter-free original-signature
term DAG with

\[
\begin{aligned}
\operatorname{size}  = O(3^r/r), \\
\operatorname{depth} = r+O(\log r).
\end{aligned}
\]

For every \(r\ge 64\), one such DAG has

\[
\begin{aligned}
\operatorname{size}  < 34\mathbin{\cdot}3^r/r, \\
\operatorname{depth} \le  r+4\mathbin{\cdot}\lceil \log_2 r \rceil+9.
\end{aligned}
\]

\textbf{Proof.} Apply Proposition 3.3 to obtain \(\sigma\), construct
the common anchor, take the union of the two DAGs in Lemmas 6.10a and
6.10b, and adjoin the glue. The nonbinary ledger already reserves the
unique glue. Therefore

\[
\begin{aligned}
\operatorname{size}/(3^r/r) < 32+1/27+1/2+1/4+1/4 < 34.
\end{aligned}
\]

Free fan-out makes this a union of ancestors, not a formula expansion:
each program vector, table root, and control root is installed once. The
depth is the maximum, not the sum, of the anchor/name path, the binary
path, and the nonbinary path, followed by two glue levels. The two
lemmas give the displayed bound. Theorem 6.1 handles the finite range
\(r<64\). \(\square\)

Together with Theorem 4.2, Theorem 6.10 matches the leading depth
coefficient one. It does not determine the optimal additive term,
ordinary-tree complexity, a bounded-fanout bound, or an exact global
size constant. A standalone no-author-import implementation
independently reconstructs the anchor, signed router, sibling-shared
vector, both compiler branches, decoder, same-DAG union, and glue; it
checks effective mutations and exact integer ledgers through arity
16,384. The integrated construction has not yet been serialized and
proved in Lean, and this internal audit is not external peer review.

\hypertarget{compiler-progression-and-bounded-validation}{%
\subsection{7. Compiler progression and bounded
validation}\label{compiler-progression-and-bounded-validation}}

\hypertarget{one-common-arity-nine-target}{%
\subsubsection{7.1 One common arity-nine
target}\label{one-common-arity-nine-target}}

The compilers were compared on the same fixed arity-nine reduction
selector, evaluating all \(3^9=19,683\) input rows. The reported node
counts include variables, and reported depths put variables at depth
one. They describe the emitted DAGs under that executable convention;
they are not \(C_r(f)\), \(D_r(f)\), exact minima, or lower bounds.

\begin{longtable}[]{@{}lrr@{}}
\toprule\noalign{}
compiler schedule & distinct reachable DAG nodes & reported depth \\
\midrule\noalign{}
\endhead
\bottomrule\noalign{}
\endlastfoot
no-wasted-leaf first-\(2\) classifier & 58,331 & 37 \\
Gray local coding & 19,329 & 179 \\
layered prefix library & 18,327 & 45 \\
recursive subcube library & 18,135 & 37 \\
global anchor and depth-three router & 23,807 & 41 \\
\end{longtable}

The table illustrates different goals. Gray local coding produces the
first large size reduction but serializes the library. Layering removes
that chain. Recursive subcubes give the best bounded result in this
test. The global-anchor compiler is asymptotically shallower but is
worse at arity nine.

The structural analytic bounds are equal at arity \(16\); the global
bound is smaller than \(4r\) from arity \(20\) onward. At arities
\(32,64,128,256\), the recorded depth-bound savings are
\(13,42,103,228\), respectively. These are arithmetic consequences of
the proved formulas, not performance claims about a materialized large
circuit.

The later routing stages are theorem-level schedule comparisons, not
additions to the materialized arity-nine benchmark above:

\begin{longtable}[]{@{}
  >{\raggedright\arraybackslash}p{(\columnwidth - 6\tabcolsep) * \real{0.2143}}
  >{\raggedleft\arraybackslash}p{(\columnwidth - 6\tabcolsep) * \real{0.2857}}
  >{\raggedleft\arraybackslash}p{(\columnwidth - 6\tabcolsep) * \real{0.2857}}
  >{\raggedright\arraybackslash}p{(\columnwidth - 6\tabcolsep) * \real{0.2143}}@{}}
\toprule\noalign{}
\begin{minipage}[b]{\linewidth}\raggedright
routing stage
\end{minipage} & \begin{minipage}[b]{\linewidth}\raggedleft
same-DAG size
\end{minipage} & \begin{minipage}[b]{\linewidth}\raggedleft
depth ceiling
\end{minipage} & \begin{minipage}[b]{\linewidth}\raggedright
evidence status
\end{minipage} \\
\midrule\noalign{}
\endhead
\bottomrule\noalign{}
\endlastfoot
historical R6 & \(O(3^r/r)\) &
\((3\mathbin{\cdot}\log_6(3))r+O(r/\log r)\) & manuscript proof plus
bounded replay and audit \\
historical Boolean-plane R9 & \(O(3^r/r)\) & \((3/2)r+O(r/\log r)\) &
manuscript proof plus bounded replay and audit \\
recursive signed family & \(O(3^r/r)\) & \(r+O(\sqrt{r})\) & exact
family independently audited and Lean-checked; historical repaired
compiler stage \\
sibling-shared vector & \(O(3^r/r)\) & \(r+O(\log r)\) & local
\((4/3+o(1))q\) vector plus matching leading scalar-output lower bound;
integrated compiler independently reconstructed \\
\end{longtable}

R9 is therefore a preserved historical stage, not the current ceiling.

\hypertarget{finite-checks}{%
\subsubsection{7.2 Finite checks}\label{finite-checks}}

The deterministic checkers currently establish the following bounded
facts.

\begin{itemize}
\tightlist
\item
  The exact category census agrees with the formula through arity eight.
\item
  All \(1\) and \(32\) semantic conservative operations at arities one
  and two are directly enumerated and compiled.
\item
  The arity-three formula evaluates to \(23,887,872\).
\item
  All \(128\) compatible arity-two selector-index tables are compiled
  and checked; the larger number reflects nonunique selector
  representations, not additional semantic operations.
\item
  Fixed-seed full-table evaluation is performed through arity seven.
\item
  Every row of the arity-nine reduction selector is evaluated.
\item
  Gray paths through block length four, layered selector censuses
  through \(M=8\), representative recursive-subcube censuses, and
  asymptotic plan arithmetic through arity \(64\) are checked.
\item
  The absorber is checked on all nine input pairs; the router is checked
  on all 81 valuations; balanced anchors are checked through arity
  seven; and global structural depth arithmetic is checked through arity
  \(256\).
\item
  The programmable routers are checked on all 4,374 six-way cases, 896
  seven-way cases, and 896 complement-relative cases. The composed
  compiler covers all 128 compatible arity-two selector tables plus
  fixed-seed tables through arity five, and its library and structural
  recurrences are checked through 128 points and arity 4,096,
  respectively.
\item
  The Boolean R9 is checked on all 4,608 projections and 9,216
  relative-orientation cases; the two-plane transfer is checked on all
  177,147 target/payload cases. Its composed compiler covers 135 tables
  and 1,641 complete tuples through arity five, with the library
  recurrence checked through 256 points and structural depth arithmetic
  through arity 4,096.
\item
  The recursive signed family is proved by induction, independently
  checked through direct truth tables and discriminator ROBDDs,
  mutation-tested on every R27 target, and Lean-checked for every depth,
  sign, mode, target, and Boolean branch valuation. The independent
  compiler audit checks every integer arity \(64..10000\) after the
  residual-first repair.
\item
  The historical E/G vector is checked on 1,200,114 target/physical/sign
  pairs, 18 dependency sets, 11,356 concrete router cases, 262
  direct-discriminator ROBDD modes, and 232 padding cases. Its author
  ledger checks every integer arity \(64..4096\); an independent audit
  extends the arithmetic replay through arity 16,384 and materializes
  both signs through width nine.
\item
  The sibling-shared vector has two independent no-author-import
  reconstructions. The broader algebraic replay checks 5,978,710
  control-pair identities across five constructions through width six,
  exact structural counts through width ten, recurrence arithmetic,
  ROBDD and concrete router semantics, and effective mutations. Its
  normal and optimized receipts are byte-identical. The separate referee
  reconstruction derives the exact recurrences and scalar-output census
  without importing the author builder.
\item
  Lean proves the optimized rail algebra, disjointness, one-digit bases,
  one-layer block composition, zero-gain reuse, sibling identity, final
  controls, exact sibling recurrences, the all-width
  \(3S\le 4q+15\mathbin{\cdot}3^{\lceil w/2 \rceil}\) envelope, and
  uniform \(S\le 7q/3\). The formal cost is a recurrence-level ledger
  rather than a serialized hash-consed DAG; the independent structural
  implementations supply that refinement evidence. Lean now proves the
  sharp local depth. It does not yet prove the complete distinct-output
  lower census or the end-to-end \(r+O(\log r)\) compiler.
\item
  An additive manuscript-repair Lean layer re-proves router
  cardinalities with kernel-reduced \passthrough{\lstinline!decide!},
  replaces vacuous signature predicates by exhaustive
  no-nullary-constructor catalogs, and proves the anchor cell, outer
  glue, and valid two-plane decoder identities. It still does not
  formalize the full ordered anchor DAG or integrated compiler.
\item
  The independent compiler-bridge reconstruction imports no author
  implementation. It checks 3,279 anchor rows, 597,870 target/physical
  vector identities, 4,632 concrete router cases, all 128 compatible
  arity-two selectors plus four deterministic arity-three selectors
  through both branches and the glue, and every exact integer ledger
  from arity 64 through 16,384. Normal and optimized outputs are
  byte-identical. Its written induction and inequalities, not that
  bounded sweep alone, support Theorem 6.10.
\item
  Mutation controls detect omitted complement quotienting, spurious
  all-\(2\) choices, complement asymmetry, corrupted Gray paths, invalid
  block caps or schedules, incorrect selector censuses, dropped depth
  corrections, wrong absorbers, constant swaps, incomplete tables, and
  binary-anchor scope errors.
\item
  Programmable-routing controls additionally reject effective gadget
  mutations, illegal absolute binary controls, omitted or per-node
  control accounting, sparse full-code accounting, omission of the
  binary depth branch, and a non-complement-invariant selector table.
\item
  Boolean-plane controls additionally reject serializing the planes,
  repeating the decoder at every level, omitting the decoder charge, and
  moving from per-level to per-node controls.
\end{itemize}

These computations validate implementations, bounded identities, and
arithmetic. They do not prove the all-\(r\) theorems or publication
novelty.

\hypertarget{replay-receipts}{%
\subsubsection{7.3 Replay receipts}\label{replay-receipts}}

Normal and optimized Python runs agree at the following semantic hashes.

\begin{longtable}[]{@{}
  >{\raggedright\arraybackslash}p{(\columnwidth - 2\tabcolsep) * \real{0.5000}}
  >{\raggedright\arraybackslash}p{(\columnwidth - 2\tabcolsep) * \real{0.5000}}@{}}
\toprule\noalign{}
\begin{minipage}[b]{\linewidth}\raggedright
receipt
\end{minipage} & \begin{minipage}[b]{\linewidth}\raggedright
semantic SHA-256
\end{minipage} \\
\midrule\noalign{}
\endhead
\bottomrule\noalign{}
\endlastfoot
\path{runs/quasiprimal_conservative_term_count/summary.json} &
\texttt{\seqsplit{f7f1758972b68abd8aba60ae547cb899c5d9376cced7cba449c34a2697e001b8}} \\
\path{runs/quasiprimal_local_coding_compiler/summary.json} &
\texttt{\seqsplit{4aaf27d678847aa0455ae8be552a165c3e9ec08c007031878f6aa40b0794ed89}} \\
\path{runs/quasiprimal_layered_local_coding_compiler/summary.json} &
\texttt{\seqsplit{b42fa56c1e78e670e2a436181337957309a0b2a658b6f539486b85cedc482d54}} \\
\path{runs/quasiprimal_subcube_local_coding_compiler/summary.json} &
\texttt{\seqsplit{f9d89ec9829c3abef05d1cdc17ef811330674bd056d90c3579a8f32b0b2eb7b8}} \\
\path{runs/quasiprimal_global_anchor_routing/summary.json} &
\texttt{\seqsplit{fb5ca7bf56d9aba9fdb0f69d1f1355c92d56d4f9306ee08a922bfc770ef60501}} \\
\path{runs/quasiprimal_programmable_projection_routing/summary.json} &
\texttt{\seqsplit{e663534756cf67ffb12fd9ae3fe722c465bd70cd91b73af1fb8084af1c0706dc}} \\
\path{runs/quasiprimal_boolean_r9_routing/summary.json} &
\texttt{\seqsplit{e4a09ef72986c1221c28e80d8227879a9e0aa0f93679e49e9f750f7e65977fce}} \\
\path{research/tournaments/2026-08-13-semantic-router-frontier/lanes/fused/summary.json}
&
\texttt{\seqsplit{4852080965ac89f54f436401950ca22ecc0dd7887593b5c01dcadcaeb883ea2a}} \\
\path{research/tournaments/2026-08-13-semantic-router-frontier/audits/strong_family/receipt.json}
&
\texttt{\seqsplit{ac566d0b716575ff693bc3a1f125283c764ef608a04c183a17293fe66adc4cb2}} \\
\path{research/tournaments/2026-08-13-semantic-router-frontier/audits/variable_compiler/receipt.json}
&
\texttt{\seqsplit{d81d0bd6b979ec0df4270fd60f04468884de6e286aea450ba28eed9897e51870}} \\
\path{research/tournaments/2026-08-13-semantic-router-frontier/lanes/program_vector/receipt.json}
&
\texttt{\seqsplit{9699a3fe058369ac22a4a4d8050d07825cc80112cbc4c601a85914100db13f88}} \\
\path{research/tournaments/2026-08-13-semantic-router-frontier/audits/program_vector/receipt.json}
&
\texttt{\seqsplit{4922df2221b3e6195d7b1af3a3ac770ae92a8aed10843430b9cffb20959c58ea}} \\
\path{research/tournaments/2026-08-13-program-vector-optimization/lanes/native_scan/receipts/receipt.json}
&
\texttt{\seqsplit{3b30a1ad6de10d39705d65edc3b27b8fd22acf361d2f7bd68f47d4d459e6d2ac}} \\
\path{research/tournaments/2026-08-13-program-vector-optimization/lanes/referee/kuhn_boolean_rail_audit/receipt.json}
&
\texttt{\seqsplit{1be842ffaca9c0510834e0eac1ee38f0acc6c310d15bb954c1a76c24dcc6f342}} \\
\path{research/tournaments/2026-08-13-program-vector-optimization/lanes/algebraic/receipt.json}
&
\texttt{\seqsplit{647f91c6c85173bbcb275e43a2301a67c61a7fd4d703b4720664a1341d7277aa}} \\
\path{research/tournaments/2026-08-13-program-vector-optimization/audits/compiler_bridge/receipt.json}
&
\texttt{\seqsplit{73a43b1e4292b56824e82a471833e69252c52f212463f182a90df63703eec23e}} \\
\end{longtable}

The formal checkpoints are source- and receipt-hash bound:

\begin{longtable}[]{@{}
  >{\raggedright\arraybackslash}p{(\columnwidth - 4\tabcolsep) * \real{0.3333}}
  >{\raggedright\arraybackslash}p{(\columnwidth - 4\tabcolsep) * \real{0.3333}}
  >{\raggedright\arraybackslash}p{(\columnwidth - 4\tabcolsep) * \real{0.3333}}@{}}
\toprule\noalign{}
\begin{minipage}[b]{\linewidth}\raggedright
formal source
\end{minipage} & \begin{minipage}[b]{\linewidth}\raggedright
source SHA-256
\end{minipage} & \begin{minipage}[b]{\linewidth}\raggedright
receipt SHA-256
\end{minipage} \\
\midrule\noalign{}
\endhead
\bottomrule\noalign{}
\endlastfoot
\path{lanes/formal/StrongSignedRouter.lean} &
\texttt{\seqsplit{687a77ed1f3b0bfe2a540bc670f6db942e15bddc339ccfbceffba9699766227a}}
&
\texttt{\seqsplit{e984f5cf50bc23489fe7bd832f08c075f61abf66885c6b0736735420eb934662}} \\
\path{lanes/formal_program_vector/ProgramVector.lean} &
\texttt{\seqsplit{c2f477edd110c4df96f3c30f31f02de09af93045720babfd38dd2dd79a5573dd}}
&
\texttt{\seqsplit{5e79e4eb49af361fbed1c8962a67583a48e532526698da584b8825d55136cffa}} \\
\path{lanes/formal_program_vector_cost/ProgramVectorCost.lean} &
\texttt{\seqsplit{ea516b085cdedd3f0ee70f83a9d0240df55e7e68cf0ad8ce77558efd91db55d2}}
&
\texttt{\seqsplit{4919835f32d8afdfc48f72979edf4ad4ef0b1f8b16271ed9858a60750d2b6211}} \\
\path{program-vector-optimization/lanes/formal_optimized/DirectRail.lean}
&
\texttt{\seqsplit{57963dadfb40cef30d94f630126f3179c55e98bea68e5974daa8e91913c404bd}}
&
\texttt{\seqsplit{cd5f1a49799726c3009748b5634a30cc154dd05d28440881df07404738c4e42e}} \\
\path{program-vector-optimization/lanes/formal_optimized/SiblingShared.lean}
&
\texttt{\seqsplit{8eb02a3e2e5a4f7531d740b704422d09724376821f83b921aec98b77a0277e77}}
&
\texttt{\seqsplit{cd5f1a49799726c3009748b5634a30cc154dd05d28440881df07404738c4e42e}} \\
\path{program-vector-optimization/lanes/formal_gap_closure/DepthAndCensus.lean}
&
\texttt{\seqsplit{ea75a228cc1316f0c53cff397409d4b9f6be4202611e29be28488a81f9a7be8c}}
&
\texttt{\seqsplit{7dcae0bbc3b3ced6d244886be01f3090a382859a6baac44906a31bb87b526a79}} \\
\path{program-vector-optimization/lanes/formal_manuscript_repairs/ManuscriptRepairs.lean}
&
\texttt{\seqsplit{4b082559f9bdcea0dc8518992c9c9a87960f2a33a289a0d3cf7ad55347b3ced7}}
&
\texttt{\seqsplit{cb9e6f0d0bbee03738e39a541e305a171e20a65fbac82b70e8f4faac59e16631}} \\
\end{longtable}

\hypertarget{relation-to-prior-work}{%
\subsection{8. Relation to prior work}\label{relation-to-prior-work}}

This draft intentionally separates classical ingredients, directly
adjacent work, and the still-unknown originality of the exact
conjunction.

\hypertarget{universal-algebra}{%
\subsubsection{8.1 Universal algebra}\label{universal-algebra}}

Pixley's discriminator and quasi-primal results {[}1,2{]} supply the
internal-isomorphism characterization used in Proposition 3.1.
Baker--Pixley interpolation {[}3{]} and Quackenbusch's extension
hierarchy {[}4{]} form the wider structural background. None of those
results is claimed here. Modern discriminator-clone and
conservative-clone work, including {[}14{]}, must also be checked in
full before submission.

\hypertarget{shannon-synthesis-and-multivalued-logic}{%
\subsubsection{8.2 Shannon synthesis and multivalued
logic}\label{shannon-synthesis-and-multivalued-logic}}

Shannon's circuit-counting method is classical {[}5{]}. Lupanov
introduced the local-coding principle used in the upper bound {[}6{]}.
Yablonskii's functional constructions in \(k\)-valued logic establish a
much broader historical setting {[}7{]}. Orlov studies realization and
basis-sensitive circuit and formula Shannon functions for \(k\)-valued
bases {[}8,27{]}. Kochergin establishes linear Shannon-depth behavior
for \(k\)-valued functions over arbitrary complete finite bases {[}9{]}.
Safin studies depth versus complexity in precomplete multivalued classes
{[}10{]}. Korshunov's survey maps local coding, Shannon effects, leading
coefficients, formulas, and depth {[}13{]}. Gashkov's primary 1978
theorem gives Boolean formulas over \(\{AND,OR,NOT\}\) with depth
\(n-\log_2(\log_2 n)+O(1)\) while retaining formula-Shannon size
\(\Theta(2^n/\log n)\) {[}28{]}.

The present parameter-free basis is incomplete because every term
preserves \(B=\{0,1\}\). Therefore the scopes are not textually
identical. That difference is not novelty evidence: the present
theorems, including the signed recurrence and sibling-shared vector, may
still be expected specializations or consequences of older closed-class
complexity, local-code decoding, switching-term, multiplexer, or
universal-circuit machinery. Dual-rail or two-rail representations are
also classical in logic synthesis; Ishiura gives a
two-rail-input/two-rail-output BDD synthesis construction in a different
basis and objective {[}19{]}, and Backes uses dual-rail encodings of
ternary semantic values for verification-oriented circuit
transformations {[}20{]}. These do not establish or refute the
fixed-\(Q\) theorem, but they prevent claiming the representation change
itself as novel. Full-text Russian and international review is required.
Tarasov's work on depth versus complexity for many-valued functional
systems is another nearby line that must be compared directly before
submission {[}17,18{]}.

The first-mismatch state law in Theorem 6.9 is also a classical
kill/propagate/generate-style prefix monoid. Ladner--Fischer
parallel-prefix computation {[}29{]} and Brent--Kung generate/propagate
networks {[}30{]} are therefore explicit methodological ancestry. The
paper does not claim the monoid or its balanced evaluation as new. The
potentially local contribution is the exact constant-free discriminator
embedding, the sibling root identities, and the matching \(4q/3\)
scalar-output census under this paper's cost model.

\hypertarget{simultaneous-size-and-depth-and-multiplexer-depth}{%
\subsubsection{8.3 Simultaneous size and depth, and multiplexer
depth}\label{simultaneous-size-and-depth-and-multiplexer-depth}}

Lozhkin gives Boolean results on formulas simultaneously approaching
Shannon size and depth {[}11{]} and, more recently, exact
basis-sensitive depth results for multiplexers with few selector lines
{[}12{]}. Kochergin's finite-complete-basis theorem {[}9{]} already
makes linear Shannon depth a classical phenomenon; the point here is one
incomplete parameter-free basis, not the order \(\Theta(r)\) by itself.
McColl and Paterson give a classical leading-one depth shape for all
Boolean functions in a richer basis {[}24{]}. Gashkov's source proof
uses a Lupanov parallel-series representation, but a checked transfer
ledger shows that dense Booleanization gives leading depth
\(\log_2(3)r\) and source-witness size \(\Theta(3^r/\log r)\), so it
does not directly supply the present same-DAG order \(O(3^r/r)\) or
coefficient one {[}28{]}. This nontransfer is not novelty evidence: a
native Q-valued Lupanov--Gashkov construction remains open, and
Lupanov's separate delayed-circuit model still requires a primary-source
model comparison. Valiant's universal circuits and Cook--Hoover
depth-universal circuits establish programmable control and joint
size/depth routing as classical themes {[}21,22{]}. Holmgren and
Rothblum give modern linear-size shared multiselection circuits in a
different Boolean model {[}23{]}. Those models do not immediately
translate to this incomplete ternary term clone, but they make broad
novelty claims about programmable controls, shared selection,
simultaneous synthesis, or leading-one depth unsafe. The unresolved
comparison is the exact conjunction of the fixed no-nullary signature,
signed P/N counts, charged sibling-shared materialization, and same-DAG
compiler. Because \(d\) has arity three on a three-element address
alphabet, a leading depth coefficient one is the natural information
rate; any eventual novelty claim must rest on the exact incomplete-basis
realization, not surprise at that coefficient.

\hypertarget{chinese-language-adjacency}{%
\subsubsection{8.4 Chinese-language
adjacency}\label{chinese-language-adjacency}}

The public-index search located Wu and Chen's 1985 ternary
sequential-circuit work with three-rail output {[}15{]} and Xu, Guan,
and Zhang's 2013 ternary reversible-logic synthesis algorithm {[}16{]}.
A focused follow-up also located Zhang, Chen, and Wang's BDD-to-MUX
area/delay optimization {[}25{]} and He, Gong, and Wei's shared
multi-output BDD construction {[}26{]}. Their algebras, reversibility
conditions, and physical, empirical, or gate-cost objectives differ from
original-signature \(Q\)-term DAGs, but the latter two are direct
adjacency for shared selector structures.

No direct match was verified in the bounded search for the exact
fixed-\(Q\) census, the d-only signed \(P_h/N_h\) family, the
sibling-shared \(4q/3\) local vector, or the integrated same-DAG
\(O(3^r/r)\) and \(r+O(\log r)\) conjunction. This is only weak negative
evidence. Native CNKI/Wanfang full-text coverage, translation variance
for clone-theoretic terminology, Russian citation-chain review,
paywalled many-valued sources, and expert comparison remain missing.

\hypertarget{paper-safe-positioning}{%
\subsubsection{8.5 Paper-safe
positioning}\label{paper-safe-positioning}}

The strongest currently supportable positioning is:

\begin{quote}
Classical interpolation, counting, local-coding, parallel-prefix, and
programmable-selection tools are specialized to one fixed parameter-free
conservative term fragment. The exact census and signed-router family
are joined by a scalar program-vector interface with matching \(4q/3\)
leading upper and lower terms in its declared local model; an explicit
two-plane and complement-relative construction then gives one same-DAG
compiler at counting-scale size and leading-one depth.
\end{quote}

Whether that exact specialization or conjunction is publication-novel
remains \textbf{UNKNOWN}.

\hypertarget{provenance-tau-boundary-and-explicit-nonclaims}{%
\subsection{9. Provenance, Tau boundary, and explicit
nonclaims}\label{provenance-tau-boundary-and-explicit-nonclaims}}

The mathematics in this draft concerns one public classical finite
algebra and public circuit-complexity methods. The recorded
implementations were independently written for OrbitSynthesis and do not
rely on Tau source, Tau specifications, private communications, private
materials, or an unsigned developer license. This is a provenance
statement only. It is not a novelty finding, a patent analysis, a
license interpretation, or a noninfringement opinion.

This draft does \textbf{not} claim:

\begin{itemize}
\tightlist
\item
  an optimal leading constant for the global Shannon-size function or
  the integrated all-arity compiler; Theorem 6.9's \(4/3\) statement is
  only the leading constant for one exact scalar-output program-vector
  interface;
\item
  an exact or optimal additive term beyond the leading-depth coefficient
  one;
\item
  that the finite Z3 result applies outside its selector-only depth-two
  grammar;
\item
  an \(\Omega(3^r)\) size lower bound;
\item
  an order-optimal ordinary unshared-term-tree bound;
\item
  novelty of Pixley interpolation, Shannon counting, Lupanov local
  coding, Gray codes, recursive Shannon expansion, layered scheduling,
  or generic linear depth;
\item
  the same count or bounds for every quasi-primal algebra or every
  incomplete clone;
\item
  a general \(r+\lceil \log_3 r \rceil+O(1)\) lower barrier for every
  compiler; the dependency argument currently applies only to a
  restricted exposed-control router cascade;
\item
  Lean extraction of a serialized materialized DAG or formal
  verification of the integrated compiler, a complete external
  independent proof review, or peer review;
\item
  publication novelty over Russian, Chinese, or international
  closed-class complexity literature;
\item
  use of, conclusions about, or limitations of Tau private work;
\item
  rights under the unsigned Tau developer license; or
\item
  patent freedom to operate.
\end{itemize}

Any commercial implementation, patent filing, categorical novelty
statement, or claim-adjacent Tau integration requires separate
legal/license review and a current prior-art search.

\hypertarget{remaining-and-quarantined-work}{%
\subsection{10. Remaining and quarantined
work}\label{remaining-and-quarantined-work}}

The mathematical next steps are:

\begin{enumerate}
\def\labelenumi{\arabic{enumi}.}
\tightlist
\item
  obtain external mathematical review of the independently reconstructed
  integrated compiler and compare its exact model with the closest
  many-valued and delayed-circuit literature;
\item
  extract or serialize the sibling-shared program-vector DAG and
  formally connect its width-indexed terms, subterm union, recurrence,
  and compiler substitution in Lean;
\item
  determine whether Lupanov's delayed-circuit model, Orlov--Lupanov,
  Kochergin, multiplexer, universal-circuit, or closed-class results
  already subsume the exact fixed-signature conjunction;
\item
  determine the optimal additive term between the counting lower bound
  and the \(r+O(\log r)\) upper bound;
\item
  determine ordinary-tree complexity and the global, rather than local
  program-vector, basis-sensitive leading-size constants; and
\item
  classify fixed finite clones for which parameter-free local coding
  gives the same order closure.
\end{enumerate}

Boolean-plane R9 and the \(7q\) E/G vector remain fully visible
historical stages. The exact signed family and sibling-shared
\((4/3+o(1))q\) vector are independently audited; the latter has a
matching leading scalar-output lower bound in its declared local model.
The complete \(r+O(\log r)\) composition has an independent
no-author-import reconstruction and is promoted as a manuscript theorem,
but not as an end-to-end Lean theorem or an externally refereed result.
Novelty, formal integration, practical value, and FTO remain
quarantined.

Negative knowledge also remains active: naive ordinary-monotone
two-extreme reasoning is false for d's antitone middle input;
residual-last growing-router chunking breaks the size ledger; sequential
\(\Theta(qw)\) program generation loses the target order; and reserving
only one logarithm can make the local library exceed the claimed bound.
Solver timeouts, bounded \passthrough{\lstinline!NO\_HIT!} results, and
grammar-specific UNSAT results are not upper bounds, lower bounds, or
optimality evidence.

\hypertarget{reproducibility-map}{%
\subsection{11. Reproducibility map}\label{reproducibility-map}}

The mathematical source notes for this version are:

\begin{itemize}
\tightlist
\item
  \path{notes/QUASIPRIMAL_CONSERVATIVE_TERM_ENUMERATION.md};
\item
  \path{notes/QUASIPRIMAL_CONSERVATIVE_TERM_SHANNON_COMPLEXITY.md};
\item
  \path{notes/QUASIPRIMAL_CONSERVATIVE_TERM_LINEAR_DEPTH.md};
\item
  \path{notes/QUASIPRIMAL_CONSERVATIVE_TERM_SUBCUBE_DEPTH.md};
\item
  \path{notes/QUASIPRIMAL_CONSERVATIVE_TERM_GLOBAL_ANCHOR_DEPTH.md};
\item
  \path{notes/QUASIPRIMAL_CONSERVATIVE_TERM_PROGRAMMABLE_DEPTH.md};
\item
  \path{notes/QUASIPRIMAL_CONSERVATIVE_TERM_BOOLEAN_R9_DEPTH.md};
\item
  \path{notes/QUASIPRIMAL_CONSERVATIVE_TERM_PARALLEL_PROGRAM_DEPTH.md};
\item
  \path{notes/QUASIPRIMAL_CONSERVATIVE_TERM_OPTIMAL_PROGRAM_VECTOR.md};
  and
\item
  \path{notes/QUASIPRIMAL_CONSERVATIVE_TERM_INTEGRATED_COMPILER.md}.
\end{itemize}

The frozen theorem, audit, and prior-art packets for the new stages are:

\begin{itemize}
\tightlist
\item
  \path{research/tournaments/2026-08-13-semantic-router-frontier/STATE.md};
\item
  \path{research/tournaments/2026-08-13-semantic-router-frontier/lanes/fused/REPORT.md};
\item
  \path{research/tournaments/2026-08-13-semantic-router-frontier/audits/strong_family/REPORT.md};
\item
  \path{research/tournaments/2026-08-13-semantic-router-frontier/audits/variable_compiler/REPORT.md};
\item
  \path{research/tournaments/2026-08-13-semantic-router-frontier/audits/variable_compiler/PRIOR_ART.md};
\item
  \path{research/tournaments/2026-08-13-semantic-router-frontier/lanes/program_vector/REPORT.md};
\item
  \path{research/tournaments/2026-08-13-semantic-router-frontier/audits/program_vector/REPORT.md};
\item
  \path{research/tournaments/2026-08-13-semantic-router-frontier/lanes/classical_depth_transfer/REPORT.md};
\item
  \path{research/tournaments/2026-08-13-semantic-router-frontier/lanes/formal/REPORT.md};
\item
  \path{research/tournaments/2026-08-13-semantic-router-frontier/lanes/formal_program_vector/REPORT.md};
\item
  \path{research/tournaments/2026-08-13-semantic-router-frontier/lanes/formal_program_vector_cost/REPORT.md};
\item
  \path{research/tournaments/2026-08-13-program-vector-optimization/lanes/native_scan/REPORT.md};
\item
  \path{research/tournaments/2026-08-13-program-vector-optimization/lanes/referee/kuhn_boolean_rail_audit/REPORT.md};
\item
  \path{research/tournaments/2026-08-13-program-vector-optimization/lanes/algebraic/REPORT.md};
\item
  \path{research/tournaments/2026-08-13-program-vector-optimization/lanes/formal_optimized/REPORT.md};
\item
  \path{research/tournaments/2026-08-13-program-vector-optimization/lanes/formal_gap_closure/REPORT.md};
\item
  \path{research/tournaments/2026-08-13-program-vector-optimization/lanes/formal_manuscript_repairs/REPORT.md};
  and
\item
  \path{research/tournaments/2026-08-13-program-vector-optimization/audits/compiler_bridge/REPORT.md}.
\end{itemize}

The implementation and deterministic checkers are:

\begin{itemize}
\tightlist
\item
  \path{src/orbitsynthesis/orbit_term_compile_local.py};
\item
  \path{experiments/quasiprimal_conservative_term_count.py};
\item
  \path{experiments/quasiprimal_local_coding_compiler.py};
\item
  \path{experiments/quasiprimal_layered_local_coding_compiler.py};
\item
  \path{experiments/quasiprimal_subcube_local_coding_compiler.py};
\item
  \path{experiments/quasiprimal_global_anchor_routing.py};
\item
  \path{experiments/quasiprimal_programmable_projection_routing.py};
\item
  \path{experiments/quasiprimal_boolean_r9_routing.py};
\item
  \path{research/tournaments/2026-08-13-semantic-router-frontier/lanes/fused/check_strong_router_family.py};
\item
  \path{research/tournaments/2026-08-13-semantic-router-frontier/audits/strong_family/audit_strong_family.py};
\item
  \path{research/tournaments/2026-08-13-semantic-router-frontier/audits/variable_compiler/check_variable_compiler.py};
\item
  \path{research/tournaments/2026-08-13-semantic-router-frontier/lanes/program_vector/check_parallel_program_vector.py};
\item
  \path{research/tournaments/2026-08-13-semantic-router-frontier/audits/program_vector/audit_program_vector.py};
\item
  \path{research/tournaments/2026-08-13-semantic-router-frontier/lanes/classical_depth_transfer/check_transfer_barriers.py};
\item
  \path{research/tournaments/2026-08-13-semantic-router-frontier/lanes/formal/StrongSignedRouter.lean};
\item
  \path{research/tournaments/2026-08-13-semantic-router-frontier/lanes/formal_program_vector/ProgramVector.lean};
\item
  \path{research/tournaments/2026-08-13-semantic-router-frontier/lanes/formal_program_vector_cost/ProgramVectorCost.lean};
\item
  \path{research/tournaments/2026-08-13-program-vector-optimization/lanes/native_scan/check_native_scan.py};
\item
  \path{research/tournaments/2026-08-13-program-vector-optimization/lanes/native_scan/check_optimal_scan.py};
\item
  \path{research/tournaments/2026-08-13-program-vector-optimization/lanes/native_scan/OptimalRailSemantics.lean};
\item
  \path{research/tournaments/2026-08-13-program-vector-optimization/lanes/referee/kuhn_boolean_rail_audit/audit_boolean_rails.py};
\item
  \path{research/tournaments/2026-08-13-program-vector-optimization/lanes/algebraic/check_algebraic_program_vector.py};
\item
  \path{research/tournaments/2026-08-13-program-vector-optimization/lanes/formal_optimized/DirectRail.lean};
\item
  \path{research/tournaments/2026-08-13-program-vector-optimization/lanes/formal_optimized/SiblingShared.lean};
\item
  \path{research/tournaments/2026-08-13-program-vector-optimization/lanes/formal_gap_closure/DepthAndCensus.lean};
\item
  \path{research/tournaments/2026-08-13-program-vector-optimization/lanes/formal_manuscript_repairs/ManuscriptRepairs.lean};
  and
\item
  \path{research/tournaments/2026-08-13-program-vector-optimization/audits/compiler_bridge/check_compiler_bridge.py}.
\end{itemize}

The semantic JSON receipts listed in Section 7 are the bounded-evidence
authority for numerical claims. The Lean checkpoint receipts bind the
formal sources and toolchain; they do not carry the asymptotic compiler
claim. The exact program-vector replay commands and frozen hashes are in
\path{lanes/program_vector/manifest.json}; the independent family
commands and hashes are in
\path{audits/strong_family/audit_manifest.json}. The source ledger is
\path{research/SOURCES.md}; the provenance/legal separation is
\path{research/IP_BOUNDARY.md}; and the current status map is
\path{research/RESEARCH_MAP_2026_08_13.md}.

Reproduction must preserve both the normal Python run and the optimized
\passthrough{\lstinline!python -O!} run and require equal semantic
hashes. A passing receipt validates only the finite scope named in that
receipt.

\clearpage

\hypertarget{declarations}{%
\subsection*{Declarations}\label{declarations}}
\addcontentsline{toc}{subsection}{Declarations}

\hypertarget{generative-ai-and-ai-assisted-technologies}{%
\subsubsection*{Generative AI and AI-assisted
technologies}\label{generative-ai-and-ai-assisted-technologies}}
\addcontentsline{toc}{subsubsection}{Generative AI and AI-assisted
technologies}

During exploratory research and preparation of this review draft, the
human project lead used OpenAI Codex and other separately invoked
large-language-model research agents to propose reformulations and proof
strategies, draft and critique mathematical arguments, generate and
review Python and Lean artifacts, and assist with organization and
copy-editing. Research Kernel and Morph were used as
research-orchestration and suggestion systems. Their outputs were
treated as untrusted suggestions: promoted claims were required to have
human-readable proofs and/or deterministic normal/optimized checkers,
independent reimplementations, mutation tests, and scoped Lean kernel
checks, as documented in Section 11 and the repository. AI systems are
not authors and cannot accept responsibility for the work. Before
submission or archival deposit, every named human author must inspect
and approve the final manuscript, verify citations and originality, and
accept full responsibility for its accuracy and integrity. The final
archival version must replace this provisional statement with a complete
retained-record inventory of the services, model or version identifiers
when available, purposes, and affected research stages. No
generative-AI-created figures are included.

\hypertarget{author-contributions-and-approval}{%
\subsubsection*{Author contributions and
approval}\label{author-contributions-and-approval}}
\addcontentsline{toc}{subsubsection}{Author contributions and approval}

Human authorship, author order, affiliations, and ORCIDs have not yet
been approved. Contribution statements and corresponding-author details
also remain pending. This review copy must not be deposited until all
named human authors approve the final text and accept accountability for
the work.

\hypertarget{data-code-and-reproducibility}{%
\subsubsection*{Data, code, and
reproducibility}\label{data-code-and-reproducibility}}
\addcontentsline{toc}{subsubsection}{Data, code, and reproducibility}

No empirical or personal data were used. Source code, proof artifacts,
deterministic checkers, and frozen receipts are available in the public
\href{https://github.com/TheDarkLightX/OrbitSynthesis}{OrbitSynthesis
repository} and in
\href{https://github.com/TheDarkLightX/OrbitSynthesis/pull/16}{draft PR
16}. Section 11 gives the exact reproduction map. The principal
manuscript source hash is
\texttt{\seqsplit{01ddf4565212f798d0adf5c4c9b16d6de9074a317a34f2a39a80051de55a5410}}.

\hypertarget{funding-and-competing-interests}{%
\subsubsection*{Funding and competing
interests}\label{funding-and-competing-interests}}
\addcontentsline{toc}{subsubsection}{Funding and competing interests}

Funding and competing-interest declarations have not yet been supplied
by the human authors and must be completed before submission. No legal
opinion about patents, freedom to operate, or license scope is made in
this manuscript.

\hypertarget{ethics-statement}{%
\subsubsection*{Ethics statement}\label{ethics-statement}}
\addcontentsline{toc}{subsubsection}{Ethics statement}

This work is mathematical and computational and involved no human
participants, animals, personal data, or clinical intervention.

\hypertarget{references}{%
\subsection*{References}\label{references}}
\addcontentsline{toc}{subsection}{References}

\begin{enumerate}
\def\labelenumi{\arabic{enumi}.}
\tightlist
\item
  A. F. Pixley, ``Functionally Complete Algebras Generating Distributive
  and Permutable Classes,'' \emph{Mathematische Zeitschrift} 114 (1970),
  361--372. \url{https://eudml.org/doc/171324}
\item
  A. F. Pixley, ``The Ternary Discriminator Function in Universal
  Algebra,'' \emph{Mathematische Annalen} 191 (1971), 167--180.
  \url{https://eudml.org/doc/162131}
\item
  K. A. Baker and A. F. Pixley, ``Polynomial Interpolation and the
  Chinese Remainder Theorem for Algebraic Systems,'' \emph{Mathematische
  Zeitschrift} 143 (1975), 165--174. \url{https://eudml.org/doc/172222}
\item
  R. W. Quackenbusch, ``Demi-Semi-Primal Algebras and Mal'cev-Type
  Conditions,'' \emph{Mathematische Zeitschrift} 122 (1971), 166--176.
  \url{https://eudml.org/doc/171592}
\item
  C. E. Shannon, ``The Synthesis of Two-Terminal Switching Circuits,''
  \emph{Bell System Technical Journal} 28 (1949).
  \url{https://www.nokia.com/bell-labs/publications-and-media/publications/the-synthesis-of-two-terminal-switching-circuits/}
\item
  O. B. Lupanov, ``On the Principle of Local Coding and the Realization
  of Functions in a Certain Class of Networks Composed of Functional
  Elements,'' \emph{Doklady Akademii Nauk SSSR} 140(2) (1961), 322--325.
  \url{https://www.mathnet.ru/eng/dan25510}
\item
  S. V. Yablonskii, ``Functional Constructions in a k-Valued Logic,''
  \emph{Trudy Matematicheskogo Instituta imeni V. A. Steklova} 51
  (1958), 5--142. \url{https://www.mathnet.ru/eng/tm1275}
\item
  V. A. Orlov, ``Complexity of Implementing Functions of k-Valued Logic
  by Circuits and Formulas in Functionally Complete Bases,''
  \emph{Discrete Applied Mathematics} 135(1--3) (2004), 223--233.
  \url{https://doi.org/10.1016/S0166-218X(02)00306-2}
\item
  A. V. Kochergin, ``On the Depth of k-Valued Logic Functions Over
  Arbitrary Bases,'' \emph{Journal of Mathematical Sciences} 233(1)
  (2018), 100--102. \url{https://doi.org/10.1007/s10958-018-3927-5}
\item
  R. F. Safin, ``On a Relation Between Depth and Complexity in
  Precomplete Classes of k-Valued Logic,'' \emph{Mathematical Problems
  of Cybernetics} 13 (2004), 223--278 {[}Russian{]}.
  \url{https://keldysh.ru/papers/2004/mvk/mvk2004_223.pdf}
\item
  S. A. Lozhkin, ``On Synthesis of Formulas Whose Complexity and Depth
  Do Not Exceed Asymptotically Best Estimates of High Accuracy,''
  \emph{Moscow University Mathematics Bulletin} 62(3) (2007), 93--100.
  \url{https://www.mathnet.ru/eng/vmumm1048}
\item
  S. A. Lozhkin, ``On the Depth of a Multiplexer Function with a Small
  Number of Select Lines,'' \emph{Mathematical Notes} 115(5) (2024),
  748--754. \url{https://www.mathnet.ru/eng/mzm14190}
\item
  A. D. Korshunov, ``Computational Complexity of Boolean Functions,''
  \emph{Russian Mathematical Surveys} 67(1) (2012), 93--165.
  \url{https://www.mathnet.ru/eng/rm9459}
\item
  E. Lehtonen and A. Szendrei, ``Equivalence of Operations with Respect
  to Discriminator Clones,'' \emph{Discrete Mathematics} 309 (2009),
  673--685. \url{https://arxiv.org/abs/0706.0195}
\item
  吴训威 and 陈偕雄,
  ``具有三轨输出的三值触发器及其在三值时序电路中的应用,'' \emph{中国科学
  A辑} 1985(7), 643--653.
  \url{https://www.sciengine.com/doi/pdf/39cfd4607c82435d895d35165b3b28c0}
\item
  徐明强, 管致锦, and 张海豹, ``基于最小混乱度的三值可逆逻辑综合算法,''
  \emph{电子学报} 41(7) (2013), 1352--1357.
  \url{https://sns.wanfangdata.com.cn/sns/perio/dianzixb/?isSync=0\&issueNum=07\&page=2\&publishYear=2013\&tabId=article}
\item
  P. B. Tarasov, ``Uniformity of Certain Systems of Functions of
  Many-Valued Logic,'' \emph{Moscow University Mathematics Bulletin}
  68(2) (2013), 99--102. \url{https://m.mathnet.ru/eng/vmumm398}
\item
  P. B. Tarasov, ``Several Conditions for Uniformity of a Finite System
  of Many-Valued Logic,'' \emph{Uchenye Zapiski Kazanskogo Universiteta.
  Seriya Fiziko-Matematicheskie Nauki} 156(3) (2014), 123--131.
  \url{https://www.mathnet.ru/eng/uzku1272}
\item
  N. Ishiura, ``Synthesis of Multilevel Logic Circuits from Binary
  Decision Diagrams,'' \emph{IEICE Transactions on Information} E76-D(9)
  (1993), 1085--1092.
  \url{https://globals.ieice.org/en_transactions/information/10.1587/e76-d_9_1085/_p}
\item
  J. Backes, \emph{Algorithms and Data Structures for Logic Synthesis
  and Verification}, Ph.D.~dissertation, University of Minnesota (2013),
  dual-rail ternary encoding discussion.
  \url{https://loonwerks.com/publications/pdf/backes2013phd.pdf}
\item
  L. G. Valiant, ``Universal Circuits (Preliminary Report),''
  \emph{Proceedings of STOC 1976}, 196--203.
  \url{https://doi.org/10.1145/800113.803649}
\item
  S. A. Cook and H. J. Hoover, ``A Depth-Universal Circuit,'' \emph{SIAM
  Journal on Computing} 14(4) (1985), 833--839.
  \url{https://doi.org/10.1137/0214058}
\item
  J. Holmgren and R. Rothblum, ``Linear-Size Boolean Circuits for
  Multiselection,'' \emph{CCC 2024}, 11:1--11:20.
  \url{https://doi.org/10.4230/LIPIcs.CCC.2024.11}
\item
  W. F. McColl and M. S. Paterson, ``The Depth of All Boolean
  Functions,'' \emph{SIAM Journal on Computing} 6(2) (1977), 373--380.
  \url{https://doi.org/10.1137/0206026}
\item
  张会红, 陈治文, and 汪鹏君, ``二叉决策图映射电路的面积和延时优化,''
  \emph{电子与信息学报} 41(3) (2019), 725--731.
  \url{https://doi.org/10.11999/JEIT180443}
\item
  何新华, 宫云战, and 魏道政, ``面向多输出电路的BDD拼接构造,''
  \emph{电子与信息学报} 19(3) (1997), 356--360.
  \url{https://jeit.ac.cn/article/id/85665992-ae08-49a4-83ca-1381a6d4735f}
\item
  V. A. Orlov, ``Realization of k-Valued Functions by Circuits of
  Functional Elements,'' \emph{Mathematical Notes} 64 (1998), 371--376.
  \url{https://doi.org/10.1007/BF02314847}
\item
  S. B. Gashkov, ``On the Depth of Boolean Functions,'' \emph{Problemy
  Kibernetiki} 34 (1978), 265--268 {[}Russian{]}.
  \url{https://publ.lib.ru/ARCHIVES/P/\%27\%27Problemy_kibernetiki\%27\%27_\%28seriya\%29/\%cf\%f0\%ee\%e1\%eb\%e5\%ec\%fb\%20\%ea\%e8\%e1\%e5\%f0\%ed\%e5\%f2\%e8\%ea\%e8.\%20\%c2\%fb\%ef\%f3\%f1\%ea\%2034.(1978).pdf}
\item
  R. E. Ladner and M. J. Fischer, ``Parallel Prefix Computation,''
  \emph{Journal of the ACM} 27(4) (1980), 831--838.
  \url{https://doi.org/10.1145/322217.322232}
\item
  R. P. Brent and H. T. Kung, ``A Regular Layout for Parallel Adders,''
  \emph{IEEE Transactions on Computers} C-31(3) (1982), 260--264.
  \url{https://doi.org/10.1109/TC.1982.1675982}
\end{enumerate}

\end{document}
